%DOCUMENTCLASS	
\documentclass[letter,11pt]{article}

%FONT ENCODINGS
\usepackage[english]{babel}
\usepackage[utf8]{inputenc}

%FONT
\renewcommand*\ttdefault{cmvtt} % change tt font to narrow

%MARGIN SPACING
\usepackage[margin=1.2 in, top=1.1 in, bottom= 1.2 in]{geometry}
\setlength {\marginparwidth }{2cm} 
 
%EQUATION SPACING
\makeatletter
\g@addto@macro\normalsize{%
  \setlength\abovedisplayskip{7pt}
  \setlength\belowdisplayskip{7pt}
  \setlength\abovedisplayshortskip{7pt}
  \setlength\belowdisplayshortskip{7pt}
}
\makeatother

%TABLE OF CONTENT LAYOUT
\usepackage{tocloft}
\renewcommand\cftsecfont{\bfseries\footnotesize}
\renewcommand\cftsubsecfont{\footnotesize}
\renewcommand\cftsecpagefont{\bfseries\footnotesize}
\renewcommand\cftsubsecpagefont{\footnotesize}
\renewcommand\cftsecafterpnum{\vskip0pt}
\renewcommand\cftsubsecafterpnum{\vskip0pt}
\setlength{\cftbeforesecskip}{0pt}
\setlength\cftparskip{0pt}
\setlength\cftaftertoctitleskip{6pt}

%FOOTNOTE LAYOUT
\interfootnotelinepenalty=10000

%ALPHABETS
\usepackage{amsfonts}
\usepackage{mathrsfs}
\usepackage{bbm}
\usepackage{latexsym}
\usepackage{math dots}
\usepackage{amssymb}
\usepackage{mathtools}
\usepackage{relsize}
\usepackage{lipsum}

%ENVIRONMENTAL PACKAGES
\usepackage{todonotes}
\usepackage{enumitem}
\setlist{nolistsep} 	%[label=(\roman{enumi}),ref=(\roman{enumi}),leftmargin=*]
\usepackage{amsthm}

%COLORS
\usepackage{xcolor}
\definecolor{Color1}{rgb}{0.0, 0.42, 0.47}%Blue
\definecolor{Color2}{rgb}{0.78, 0.11, 0.0}%Scarlet

%SECTION TITLE LAYOUT 
\usepackage{titlesec}
\titleformat{\section}
  {\Large\center\bfseries}
  {\thesection.}{.7em}{}
\titlespacing*{\section}{0pt}{3.5ex plus 0ex minus 0ex}{1.5ex plus 0ex}
\titleformat{\subsection}
  {\center\bfseries}
  {\thesubsection.}{.7em}{}
\titlespacing*{\subsection}{0pt}{3.5ex plus 0ex minus 0ex}{1.5ex plus 0ex}
\titleformat{\subsubsection}
  {\center\bfseries}
  {\thesubsubsection.}{.7em}{}
\titlespacing*{\subsubsection}{0pt}{3.5ex plus 0ex minus 0ex}{1.5ex plus 0ex}

%CONTENTS TITLE LAYOUT
\addto\captionsenglish{\renewcommand{\contentsname}{\large\bfseries Contents}}

%DOCUMENT TITLE LAYOUT
\usepackage{titling}
\setlength{\droptitle}{-4em}
%\posttitle{\par\end{center}\vspace{-.7em}}

%HYPERREF PACKAGES
\usepackage[linktocpage=true]{hyperref}
\usepackage[capitalize]{cleveref}
\hypersetup{citecolor = Color1,colorlinks,
			linkcolor = black,
			urlcolor = Color2}

%AMSTHM SETUP
\newtheoremstyle{plain}{3mm}{3mm}{\slshape}{}{\bfseries}{.}{.5em}{}
\newtheoremstyle{definition}{2mm}{2mm}{}{}{\bfseries}{.}{.5em}{}
\theoremstyle{plain} %PLAIN
\newtheorem*{thm}{Theorem}	
\newtheorem{theorem}{Theorem}[section]
\newtheorem{problem}[theorem]{Problem}
\newtheorem{proposition}[theorem]{Proposition}
\newtheorem{question}[theorem]{Question}
\newtheorem{conjecture}[theorem]{Conjecture}
\newtheorem{lemma}[theorem]{Lemma}
\newtheorem{corollary}[theorem]{Corollary}
\theoremstyle{definition} %DEFINITION
\newtheorem{defn}[theorem]{Definition}
\newtheorem{remark}[theorem]{Remark}
\newtheorem{example}[theorem]{Example}

%MAIN THEOREMS
\theoremstyle{plain} 
\newcounter{MainTheoremCounter}
\renewcommand*{\theMainTheoremCounter}{\Alph{MainTheoremCounter}}
\newtheorem{maintheorem}[MainTheoremCounter]{Theorem}

%MAIN COROLLARIES
\theoremstyle{plain} 
\newcounter{MainConjectureCounter}
\renewcommand*{\theMainConjectureCounter}{\Alph{MainConjectureCounter}}
\newtheorem{mainconjecture}[MainConjectureCounter]{Conjecture}

%NAMED THEOREMS
\theoremstyle{plain}
\newtheorem*{namedthm}{\namedthmname}
\newcounter{namedthm}
\makeatletter
	\newenvironment{named}[2]
	{\def\namedthmname{#1}
	\refstepcounter{namedthm}
	\namedthm[#2]\def\@currentlabel{#1}}
	{\endnamedthm}
\makeatother
	%Example:
	%\begin{named}{Incompleteness Theorem}{\cite{Godel1931}}
	%\label{theorem:incompleteness}
	% satement of theorem
	%\end{named}
	%In 1931 Kurt G{\"o}del proved the \ref{theorem:incopleteness}.

%APPENDIX LAYOUT
\usepackage{chngcntr}
%\counterwithin{theorem}{section}

%NUMBERWITHIN
\numberwithin{equation}{section}

%ALLOWDISPLAYBREAKS
\allowdisplaybreaks

%PEEPS
\newcommand{\Caratheodory}{Carath\'{e}odory}
\newcommand{\Cech}{\v{C}ech}
\newcommand{\Cesaro}{Ces\`{a}ro}
\newcommand{\Erdos}{Erd\H{o}s}
\newcommand{\Folner}{F\o{}lner}
\newcommand{\Poincare}{Poincar\'{e}}
\newcommand{\Szemeredi}{Szemer\'{e}di}
\newcommand{\Beiglbock}{Beiglb\"{o}ck}
\newcommand{\Fejer}{Fej\'{e}r}
\newcommand{\Sarkozy}{S\'{a}rk\"{o}zy}
\newcommand{\Katai}{K\'{a}tai}
\newcommand{\Mobius}{M\"{o}bius}
\newcommand{\KulagaPrzymus}{Ku{\l}aga-Przymus}
\newcommand{\Lemanczyk}{Lema\'{n}czyk}
\newcommand{\Gauss}{Gau{\ss}}
\newcommand{\Turan}{Tur{\'a}n}
\newcommand{\Matomaki}{Matom\"{a}ki}
\newcommand{\Radziwill}{Radziwi\l\l{}}
\newcommand{\Halasz}{Hal\'{a}sz}

%STANDARD COMMANDS
\newcommand{\Oh}{{\mathrm O}}
\newcommand{\oh}{{\mathrm o}}
\newcommand{\N}{\mathbb{N}}
\newcommand{\Z}{\mathbb{Z}}
\newcommand{\R}{\mathbb{R}}
\newcommand{\C}{\mathbb{C}}
\newcommand{\Q}{\mathbb{Q}}
\newcommand{\T}{\mathbb{T}}
\newcommand{\U}{\mathbb{U}}
\newcommand{\Hilb}{\mathscr{H}}
\newcommand{\supp}{\operatorname{supp}}
%\newcommand{\Aut}{\mathrm{Aut}}
%\newcommand{\Cont}{\mathrm{C}}
%\newcommand{\Lip}{\mathrm{Lip}}
%\newcommand{\Hom}{\mathrm{Hom}}

%LAYOUT COMMANDS
\newcommand{\define}[1]{{\itshape #1}}

%RENEWED COMMANDS
\renewcommand{\epsilon}{\varepsilon}
\renewcommand{\leq}{\leqslant}
\renewcommand{\geq}{\geqslant}
\renewcommand{\setminus}{\backslash}
\renewcommand{\Re}{\mathrm{Re}}
\renewcommand{\Im}{\mathrm{Im}}
\renewcommand{\P}{\mathbb{P}}
\renewcommand{\subset}{\subseteq}
\renewcommand{\supset}{\supseteq}
\renewcommand{\d}{~\mathrm{d}}

%MULTIPLICATIVE NUMBER THEORY COMMANDS
\newcommand{\mob}{\boldsymbol{\mu}}
\newcommand{\lio}{\boldsymbol{\lambda}}
\newcommand{\tot}{\boldsymbol{\varphi}}

%EDITING COMMANDS
\newcommand{\edit}[3]{\color{#1}{#3}\color{black}\marginpar{\textcolor{#1}{[[#2]]}}}
\newcommand{\dc}[1]{\edit{blue}{DC}{#1}}
\newcommand{\fkr}[1]{\edit{red}{FR}{#1}}
\newcommand{\new}[1]{\edit{blue}{NEW}{#1}}
\newcommand{\rem}[1]{\edit{purple}{REMARK}{#1}}
\usepackage[normalem]{ulem}
\newcommand{\rsout}[1]{{\color{red}\sout{#1}}}
\usepackage[makeroom]{cancel}
\newcommand{\rcancel}[1]{{\color{red}\xcancel{#1}}}

%DOCUMENT SPECIFIC COMMANDS
\newcommand{\F}{\mathcal{F}}
\newcommand{\J}{\mathcal{J}}
\newcommand{\I}{\mathcal{I}}
\newcommand{\A}{\mathcal{A}}
\newcommand{\AN}{\mathcal{A}_\N}
\newcommand{\M}{\mathcal{M}}
\newcommand{\D}{\mathcal{D}}
\newcommand{\B}{\mathcal{B}}

\newcommand{\E}{\raisebox{-0.4ex}{\scalebox{1.4}[1.3]{$\mathbb{E}$}}}
\newcommand{\Esub}[1]{\raisebox{-0.4ex}{\scalebox{1.4}[1.3]{$\mathbb{E}$}}\raisebox{-0.35ex}{\ensuremath{_{\scriptstyle #1}}}}
\newcommand{\logE}{{\E}^{\log}}

\newcommand{\Inv}{\mathcal{I}}
\newcommand{\Dist}{\mathbb{D}}
\newcommand{\1}{\mathbbm{1}}
\newcommand{\Lip}{{\texttt{\textup{Lip}}}}
\newcommand{\mmod}{\texttt{\textup{mod}}}


%OTHER COMMANDS
\newcommand{\f}{f_{\xi,N}}
\newcommand{\Pc}{\mathcal{P}}

%ADDING REFS IN CONTENTS
\usepackage[nottoc]{tocbibind}


%==========================================================
%==========================================================

%TITLE&ABSTRACT
\author{By~~{\scshape Dimitrios~Charamaras}~~and~~{\scshape Florian~K.~Richter}}
\date{\small \today}
\title{\bfseries 
Asymptotic independence of $\Omega(n)$ and $\Omega(n+1)$ along logarithmic averages}


%BEGIN DOCUMENT
\begin{document}


\maketitle


\begin{abstract}
Let $\Omega(n)$ denote the number of prime factors of a positive integer $n$ counted with multiplicities. We show that for any bounded functions $a,b\colon\N\to\C$,
$$\frac{1}{\log{N}}\sum_{n=1}^N \frac{a(\Omega(n))b(\Omega(n+1))}{n}
= \Bigg(\frac{1}{N}\sum_{n=1}^N a(\Omega(n))\Bigg)\Bigg(\frac{1}{N}\sum_{n=1}^N b(\Omega(n))\Bigg) + \oh_{N\to\infty}(1).$$
This generalizes a theorem of Tao on the logarithmically averaged two-point correlation Chowla conjecture.
Our result is quantitative and the explicit error term that we obtain establishes double-logarithmic savings.
As an application, we obtain new results about the distribution of $\Omega(p+1)$ as $p$ ranges over $\ell$-almost primes for a ``typical'' value of $\ell$.
\end{abstract}



\tableofcontents
\thispagestyle{empty}

%==========================================================
%==========================================================


%SECTION
\section{Introduction}\label{Intro}

For a positive integer $n\in\N=\{1,2,3,\ldots\}$, let $\Omega(n)$ denote the number of prime factors of~$n$ counted with multiplicities.
Given $N\in\N$ we use the notation $[N]$ for the set of integers $\{1,\ldots,N\}$, we denote the \Cesaro{} average of a function $f\colon\N\to\C$ over a finite non-empty set $A\subset\N$ as
$$\Esub{n\in A}f(n) = \frac{1}{|A|}\sum_{n\in A}f(n),$$
and the logarithmic average of $f$ over $A$ as 
$$\logE_{n\in A}f(n)=\frac{1}{\sum_{n\in A}\frac{1}{n}}\sum_{n\in A} \frac{f(n)}{n}.$$

%The distribution of the values of $\Omega(n)$ has been studied extensively.
%One of the earliest results in this direction belongs to von Mangoldt \cite[p.\ 852]{vonMangoldt97}, who showed that the asymptotic density of~$n$ for which $\Omega(n)$ is even equals $\frac{1}{2}$, a statement that Landau later understood to be elementary equivalent to the Prime Number Theorem (see \cite{Landau11,Landau12} and \cite[pp.\ 571, 620--621, 631--632]{Landau09b}).
%Hardy and Ramanujan \cite{HR17} showed that the \define{normal order} of~$\Omega(n)$ equals $\log\log(n)$, which means that for every $\epsilon>0$ the asymptotic density of the set $\{n\in\N: (1-\epsilon)\log\log(n)\leq \Omega(n)\leq (1+\epsilon)\log\log(n)\}$ equals~$1$.
%In 1940, \Erdos{} and Kac \cite{EK40} extended the Hardy-Ramanujan Theorem by proving that, for $N\to\infty$, the distribution of $\Omega(n)$ as $n$ ranges from $1$ to $N$ approximates the distribution of a Gaussian random variable with mean value $\log\log(N)$ and standard deviation~$\sqrt{\log\log(N)}$. Their result is often considered a ``Central Limit Theorem'' for $\Omega(n)$.
%
%On smaller scales, it is widely believed that the distribution of $\Omega(n)$ does not follow any deterministic patterns.
%In particular, there appears to be no significant correlation between the values of $\Omega(n)$ along consecutive integers $n, n+1, \ldots, n+k-1$. This relates to a famous conjecture posed by Chowla concerning the Liouville function $\lio(n)\coloneqq (-1)^{\Omega(n)}$.
The starting point of our discussion is an old conjecture of Chowla concerning the Liouville function $\lio(n)\coloneqq (-1)^{\Omega(n)}$.
It asserts that asymptotically there is no correlation among the sign patterns of $\lio(n)$ along consecutive integers.


\begin{named}{Chowla's Conjecture}{\cite{Chowla87}; see also \cite{Ramare18}}
\label{conj_chowla}
For any $k\in\N$,
\begin{equation}
\label{eqn_chowla}
\Esub{n\in[N]} \lio(n) \lio(n+1)\cdots \lio(n+k-1) = \oh_{N\to\infty}(1).
\end{equation}
\end{named}
%Chowla's Conjecture implies that the sign patterns of the Liouville function, which encode the parity of the values of $\Omega(n)$, behave independently along consecutive integers.

If the \Cesaro{} average in \eqref{eqn_chowla} is replaced by a logarithmic average, then the case $k=2$ of \eqref{eqn_chowla} was proved in \cite{Tao16}.

\begin{theorem}[{cf.~\cite[Theorem 2]{Tao16}}]\label{thm_log_2_correlation_chowla}
We have
\[
\logE_{n\in[N]} \lio(n)\lio(n+1)= \oh_{N\to\infty}(1).
\]
\end{theorem}

In \cite{helf_radz}, Helfgott and \Radziwill{} established double-logarithmic savings in the above, while the best known result to date is due to Pilatte \cite{pilatte2023improved} who proved \cref{thm_log_2_correlation_chowla} with logarithmic savings.

The method used in \cite{Tao16} was later extended in \cite{TT18} to establish the logarithmically averaged \ref{conj_chowla} for any odd  $k$. 
Albeit significant recent progress combining ideas from number theory, higher order Fourier analysis, and ergodic theory (cf.\ \cite{Sarnak11,EAKLdlR17,Frantzikinakis17,tao2017equivalence,FH18a,FKL18,GLdlR21,KMT23_almost_all_scales,MRTT2023_higher_uniformity,MSTT2023,KLRT2024,MRSTT2024}),
Chowla's conjecture, and the more general Elliott's conjecture\footnote{The original Elliott's conjecture stated by Elliott in \cite{Elliott92} is false (see \cite[Theorem B.1]{MRT15} for the counterexample). The corrected version of the conjecture was given in \cite{MRT15} and it is still open.} (cf.~\cite{Elliott92,MRT15}), remain open for other $k$, even in the case of logarithmic averages.


% We remark that if Chowla's Conjecture was known to hold for $k=1$ then \cref{thm_local_2-fourier_uniformity_lio} would follow from it by the Cauchy-Schwarz inequality with an optimal error in terms of $H$ of the form $\Oh(H^{-1})+\oh_{N\to\infty}(1)$. \fkr{[[I think it should be $\Oh(H^{-\frac{1}{2}})$ instead of $\Oh(H^{-1})$, please double-check. ]]}

A different relation for the sequence $\Omega(n)$, introduced in \cite{Richter_PNT_21}, draws inspiration from an ergodic-theoretic perspective.

\begin{theorem}[cf.~{\cite{Richter_PNT_21}}]
\label{thm_Omega_is_good_for_UE}
For any bounded arithmetic function $a\colon\N\to\C$, we have
\begin{equation}
\label{eqn_pnt_1}
\Esub{n\in[N]} a(\Omega(n)) = \Esub{n\in[N]} a(\Omega(n)+1) +\oh_{N\to\infty}(1).
\end{equation}
\end{theorem}
We remark that if $a(n)=(-1)^n$ then \eqref{eqn_pnt_1} reduces to the $k=1$ case of Chowla's conjecture, %coincides with the classical Liouville function $\lio(n)=(-1)^{\Omega(n)}$ and \eqref{eqn_pnt_1} is equivalent to $\frac{1}{N}\sum_{n=1}^N\lio(n) =\oh_{N\to\infty}(1)$,
which is a well-known equivalent form of the Prime Number Theorem.
For generalizations of \eqref{eqn_pnt_1} and its connections to ergodic theory see~\cite{BR22,Loyd23, Charamaras_MVT_24}. 

Theorems~\ref{thm_log_2_correlation_chowla}~and~\ref{thm_Omega_is_good_for_UE} highlight different number-theoretic aspects of the distribution of the values of $\Omega(n)$ and 
our goal is to unify these approaches.
% and opening up unexplored directions of research.
Instead of restricting to the joint distribution of multiplicative parity values, as in Chowla's conjecture, it is natural to inquire about more general forms of independence regarding the joint distribution of $\Omega(n),\Omega(n+1),\ldots,\Omega(n+k-1)$. 
Based on \cref{thm_Omega_is_good_for_UE}, we propose the following conjecture. %, which widens the scope whereby we measure the multiplicative independence of consecutive numbers.

% given our understanding of the behaviour of $\Omega(n)$, it is reasonable to expect that if $H$ and $N$ are two large quantities, where $N$ is larger than $H$, then the average of $a(\Omega(n))$ over the ``long'' interval $\{1,2,\ldots,N\}$ approximately equals the average of $a(\Omega(n))$ over the ``short'' interval $\{n+1, n+2, \ldots, n+H\}$ for almost all starting points $n\in\{1,\ldots,N\}$.

\begin{conjecture}
\label{conj_functional_chowla}
For any $k\in\N$ and any bounded functions $a_1,\ldots,a_k\colon\N\to\C$,
$$
\Esub{n\in[N]} a_1(\Omega(n))a_2(\Omega(n+1))\cdots a_k(\Omega(n+k-1))
= \prod_{i=1}^k \Esub{n_i\in[N]} a_i(\Omega(n_i)) + \oh_{N\to\infty}(1).
$$
\end{conjecture}

Notice that if $a_1(n)=\ldots=a_k(n)=(-1)^{n}$, then \cref{conj_functional_chowla} reduces to Chowla's conjecture. More generally, if $a_1,\ldots,a_k\colon\N\to\mathbb{S}^1=\{z\in\C: |z|=1\}$ are group characters of $(\N,+)$, then \cref{conj_functional_chowla} becomes a special case of Elliott's conjecture. The intended upgrade of \cref{conj_functional_chowla} is that no structural restrictions, particularly no algebraic ones, are imposed on the functions $a_1,\ldots,a_k$.

\subsection{Main results}

The aim of this paper is to validate the logarithmically averaged two-point correlation case of \cref{conj_functional_chowla}. %\new{For simplicity, we focus on the case $(h_1,h_2)=(0,1)$.} 
Our main result is as follows:

\begin{maintheorem}\label{MT}\label{thm_A}
For any bounded functions $a,b\colon\N\to\C$, we have
    \begin{equation}\label{MT_eq}
        \logE_{n\in[N]} a(\Omega(n))b(\Omega(n+1)) 
        = \Big(\Esub{n\in[N]}a(\Omega(n))\Big)\Big(\Esub{n\in[N]}b(\Omega(n)\Big)
        + \Oh\Bigg(\frac{\sqrt{\log\log\log{N}}}{(\log\log{N})^\frac{1}{6}}\Bigg)
    \end{equation}
    Moreover, the error term only depends on $\max\{\|a\|_\infty,\|b\|_\infty\}$.
\end{maintheorem}

\cref{thm_A} is a generalization of \cref{thm_log_2_correlation_chowla}, since choosing $a(n)=b(n)=(-1)^n$ in \eqref{MT_eq} implies \eqref{eqn_chowla}.

We remark that the best bound in \eqref{MT_eq} that one could hope for is 
\[
\Oh\bigg(\frac{1}{\sqrt{\log\log{N}}}\bigg),
\]
since this bound is known to be best possible for the \Erdos{}-Kac theorem due to the work of \cite{renyi-turan}, and \eqref{MT_eq} contains the (logarithmically averaged) \Erdos{}-Kac theorem as a special case.
% this error term arises from the error term in Tao's identity \dc{[REF]}. We note that a more careful proof of this result can yield a slightly improved error term with $\frac{1}{3}$ replaced by $\frac{1}{2}$, which is not much of an improvement though. 
% The optimal bound for \eqref{MT_eq} is 
% $$\Oh\Bigg(\frac{1}{\sqrt{\log\log{N}}}\Bigg),$$
% and to prove \cref{MT} with this error term, one has to establish \dc{[REF]} with the same error term. 
% With the current input coming from the work of Helfgott and \Radziwill{} \cite{helf_radz}, it would be possible to establish \eqref{MT_eq} for $a(n,m) = a_1(n)a_2(m)$ and $h=1$. \dc{Maybe $h=1$ is not so important, and it is possible to use the Helfgott-\Radziwill{} theorem for any $h\in\N$.}

\cref{thm_A} demonstrates a natural form of asymptotic independence between $\Omega(n)$ and $\Omega(n+1)$.
We can provide a clearer, more intuitive, understanding of this behavior by recasting our result in the language of probability theory, in the spirit of the \Erdos{}-Kac theorem.
Recall that if $X$ is a Gaussian random variable with mean $\mu=\mathbb{E}[X]$ and standard deviation $\sigma=\sqrt{\mathbb{E}[(X-\mu)^2]}$ then the probability density function of $X$ is given by
\begin{equation}
\label{eqn_Gaussian_prob_dens_fctn}
f(x;\mu,\sigma)= \frac{1}{\sigma\sqrt{2\pi}}\,e^{-\frac{1}{2}\left(\frac{x-\mu}{\sigma}\right)^2}.
\end{equation}
This means that for any bounded and continuous $g\colon \R\to \C$ we have
\begin{equation}
\label{eqn_Gaussian_density_function_characterization}
\mathbb{E}[g(X)] = \int_{-\infty}^\infty g(x)f(x;\mu,\sigma) \d x. 
\end{equation}
In analogy, the ``mean'' and ``standard deviation'' of $\Omega(n)$ over $[N]$ are defined as
\begin{equation}\label{mu_N and sigma_N}
    \mu_N = \log\log{N} \sim \Esub{n\in[N]}\Omega(n)
    \quad\text{and}\quad
     \sigma_N = \sqrt{\log\log{N}} \sim \Big(\Esub{n\in[N]}\big|\Omega(n)-\mu_N\big|^2\Big)^\frac{1}{2}.
\end{equation}

Given $A>1$, we define the {\em typical range} $T_{A,N}$ of $\Omega(n)$ as
\begin{equation}\label{typical_range}
    T_{A,N} = [\mu_N-A\sigma_N,\mu_N+A\sigma_N] = [\log\log{N}-A\sqrt{\log\log{N}}, \log\log{N}+A\sqrt{\log\log{N}}].
\end{equation}
% \mu_N= \frac{1}{N}\sum_{n=1}^N\Omega(n)\sim \sum_{p\leq N} \frac{1}{p}\quad\text{and}\quad
% \sigma_N=  \frac{1}{N}\sum_{n=1}^N\big|\Omega(n)-\mu_N\big|^2\sim\bigg(\sum_{p\leq n} \bigg(\frac{1}{p}-\frac{1}{p^2}\bigg)\bigg)^{1/
% 2},
% \]
%where the sums are taken over all primes smaller or equal than $N$.
Let $\P_\ell=\{n\in\N: \Omega(n)=\ell\}$ denote the set of all $\ell$-almost primes.
Based on the uniform asymptotic estimates of $\P_\ell\cap[N]$ obtained by \Erdos{} in \cite{Erdos48a} (see also \cite{Selberge_Sathe}), one can show (see \cite[Lemma 3.4]{Loyd23} and \cref{pi_k_estimate} below) that for all $A>1$ and uniformly over all $\ell$ in the typical range $T_{A,N}$, one has
\begin{equation*}
\label{eqn_loy_lemma3.4}
\Esub{n\in[N]} \1_{\P_\ell}(n) = 
f(\ell;\mu_N,\sigma_N)\bigg(1 + \Oh\bigg(\frac{1}{\sqrt{\log\log{N}}}\bigg)\bigg).
\end{equation*}
Combined with the Hardy-Ramanujan theorem, this proves that uniformly over all functions $a\colon\N\to\C$ with $\|a\|_\infty\leq 1$,
\begin{align}
\label{eqn_loy_lemma3.4+}
\Esub{n\in[N]} a(\Omega(n))&=\sum_{\ell=1}^{\infty} a(\ell)f(\ell; \mu_N, \sigma_N)+\oh_{N\to\infty}(1).
\end{align}
Juxtaposing \eqref{eqn_Gaussian_density_function_characterization} and \eqref{eqn_loy_lemma3.4+}, it becomes clear that $f(\ell; \mu_N, \sigma_N)$ can be considered the ``probability density function'' of $\Omega(n)$ over the interval $[N]$.
Note that \eqref{eqn_loy_lemma3.4+} is a strengthening of the \Erdos{}-Kac theorem and contains \eqref{eqn_pnt_1} as a special case. It also immediately implies the following reformulation of \cref{MT}.

\begin{maintheorem}
\label{thm_B}
For any bounded functions $a,b\colon\N\to\C$, we have
\begin{equation*}
    \logE_{n\in[N]} a(\Omega(n))b(\Omega(n+1))
    =
    \sum_{k,\ell=1}^{\infty} a(k)b(\ell)f(k; \mu_N, \sigma_N) f(\ell; \mu_N, \sigma_N)
    + \Oh\Bigg(\frac{\sqrt{\log\log\log{N}}}{(\log\log{N})^\frac{1}{6}}\Bigg).
\end{equation*}
where $f(x;\mu,\sigma)$ is the probability density function of the Gaussian distribution defined in \eqref{eqn_Gaussian_prob_dens_fctn}.
Moreover, the error term only depends on $\max\{\|a\|_\infty,\|b\|_\infty\}$.
\end{maintheorem}

As an application of our main theorem, we obtain new results related to the asymptotic behavior of shifted almost primes. In particular, the next result asserts that for any bounded function $a\colon\N\to\C$, and for ``almost all'' $\ell$ in the typical range $T_{A,N}$, the sequence $a(\Omega(n+1))$ with $n$ ranging over $\ell$-almost primes asymptotically behaves the same as the sequence $a(\Omega(n))$ with $n$ ranging over all positive integers.
To state our result in this direction, we define $\overline{\pi}_\ell(N)$ to be the relative density of the set of $\ell$-almost primes within the set of integers up to $N$, that is,
\begin{equation}\label{pi_defn}
    \overline{\pi}_\ell(N)=\frac{|\P_\ell\cap [N]|}{N}.
\end{equation}

\begin{maintheorem}
\label{thm_C}
Let $a\colon\N\to\C$ be a bounded function. Then we have
\begin{equation}\label{eqn1_thm_C}
    \sum_{\ell=1}^\infty \overline{\pi}_\ell(N)\Big|\logE_{n\in[N]\cap\P_\ell} a(\Omega(n+1)) - \Esub{n\in[N]} a(\Omega(n))\Big|
    \ll_{\|a\|_\infty} \frac{\sqrt{\log\log\log{N}}}{(\log\log{N})^\frac{1}{6}}.
\end{equation}
In particular, for any fixed $A>1$ and any $\epsilon>0$, for all but at most $\epsilon|T_{A,N}|$ numbers $\ell\in T_{A,N}$, we have 
\begin{equation}\label{eqn2_thm_C}
    \logE_{n\in[N]\cap\P_\ell} a(\Omega(n+1)) = \Esub{n\in[N]} a(\Omega(n)) + \Oh_{\epsilon,\|a\|_\infty}\Bigg(\frac{\sqrt{\log\log\log{N}}}{(\log\log{N})^\frac{1}{6}}\Bigg).
\end{equation}
\end{maintheorem}

We remark that if $a(n)=(-1)^n$, then \eqref{eqn1_thm_C} with better bounds than we have is given by \cite[Eq.~(1.16)]{helf_radz}.

We prove that Theorems~\ref{MT} and~\ref{thm_C} are equivalent in \cref{Proof_Thm_C}.


\subsection{Proof strategy for \cref{thm_A}}

In broad strokes, we follow the approach introduced by Tao in \cite{Tao16} to prove \cref{thm_log_2_correlation_chowla}, which builds on results of \Matomaki-\Radziwill-Tao in \cite{MRT15}.
The biggest difference in our approach is that the results in \cite{MRT15} cannot be used directly, as they apply exclusively to \emph{strongly aperiodic}\footnote{A multiplicative function $f\colon\N\to\C$ is called \emph{strongly aperiodic} if for any $B>0$, we have $\lim_{N\to \infty}\inf_{|t|\leq BX,q\leq B}\Dist(f,n\mapsto\chi(n) n^{it};N) = \infty$,
where the infimum is taken over all Dirichlet characters $\chi$ of conductor
$q$ with $q\leq B$, and $\Dist$ denotes the \emph{distance} between two multiplicative functions (see \cref{defn_distance}).} multiplicative functions, and the functions we are dealing with here are not even multiplicative.

The starting point in the proof of \cref{MT} is reducing our problem to the case of multiplicative functions. After replacing the \Cesaro{} averages in the left hand side of \eqref{MT_eq} with logarithmic ones, then using Fourier analysis, the \Erdos-Kac theorem and a theorem of Helfgott and \Radziwill{} in \cite{helf_radz}, we can reduce bounding 
$$\Big|\logE_{n\in[N]}a(\Omega(n))b(\Omega(n+1)) - \Big(\Esub{n\in[N]}a(\Omega(n))\Big)\Big(\Esub{n\in[N]}b(\Omega(n))\Big)\Big|$$
to bounding
$$\sum_{\xi\in I_N}\logE_{n\in[N]}\Big|\logE_{p\in\Pc_N} \f(n+p) - \logE_{m\in[N]}\f(m)\Big|^2,$$
where $(\f)_{\xi\in I_N}$ is a sequence of completely multiplicative functions defined in \eqref{defn_f}, $I_N$ is a finite integer interval of frequencies $\xi$ defined in \eqref{I_N_def}, the length of which grows with $N$ in appropriate speed, and $\Pc_N$ is a set of primes with size depending on $N$ defined in \eqref{defn_P}. This reduction is carried out in \cref{reduction_section}.

Now, if $\f$ was strongly aperiodic for all frequencies $\xi\in I_N$, then one could use the results from \cite{MRT15} (after using a circle method estimate) to bound the expression in question. Note that this would not be enough though, since the sum over $\xi$ could lead the error term to explode, thus, one would need a quantitative lower bound on $\inf_{|t|\leq BX,q\leq B}\Dist(\f,n\mapsto\chi(n) n^{it};N)$, and not just divergence to infinity.
The major issue here is that not all $\f$ are strongly aperiodic. To overcome this, we distinguish between the case when the frequency $\xi$ is small and the case when it is large. For small frequencies, we bound the expression in question by establishing quantitative novel bounds for a local \Erdos-Kac theorem. This is carried out in \cref{small_xi_section}. Finally, for the large frequencies $\xi$, the function $\f$ is not just strongly aperiodic, but we have convenient lower bound for $\inf_{|t|\leq BX,q\leq B}\Dist(\f,n\mapsto\chi(n) n^{it};N)$. Using a circle method estimate and a result from \cite{MRT15}, we conclude the proof. This is carried out in \cref{large_xi_section}.


\subsection{Conjectures concerning almost primes}

Establishing, for $h,\ell\in\N$, the asymptotic growth of expressions of the form
\[
\sum_{n=1}^N \1_{\P_\ell}(n)\1_{\P_\ell}(n+h)
\]
is a longstanding and very difficult problem in number theory.
Based on the heuristic that the prime factorization of $n$ and $n+h$ display a high degree of independence, it is reasonable to expect that for all $A>0$, all $k\in\N$ and all $\ell_1,\ldots,\ell_k$ in the typical range $T_{A,N}$, we have
\begin{align}
\label{eqn_almost_primes_independence}
\Esub{n\in[N]} \1_{\P_{\ell_1}}(n)\1_{\P_{\ell_2}}(n+1)\cdots \1_{\P_{\ell_k}}(n+k-1)
&= \overline{\pi}_{\ell_1}(N)\cdots \overline{\pi}_{\ell_k}(N)(1+\oh_{N\to\infty}(1)).
\end{align}
The following conjecture implies that \eqref{eqn_almost_primes_independence} holds for ``almost all'' $\ell_1,\ldots,\ell_k$ in the typical range.

\begin{conjecture}
\label{conj_independence_almost_primes}
For all $k\in\N$, we have
\begin{equation*}
\sum_{\ell_1,\ldots,\ell_k=1}^\infty \Big|\Esub{n\in[N]} \1_{\P_{\ell_1}}(n)\1_{\P_{\ell_2}}(n+1)\cdots \1_{\P_{\ell_k}}(n+k-1) - \overline{\pi}_{\ell_1}(N)\cdots\overline{\pi}_{\ell_k}(N)\Big|
= \oh_{N\to\infty}(1).
\end{equation*}
\end{conjecture}

\cref{conj_independence_almost_primes} implies \cref{conj_functional_chowla}, although the relationship between the two is not immediately apparent. To clarify the connection, we offer below an equivalent formulation of \cref{conj_independence_almost_primes}.

\begin{conjecture}
\label{conj_functional_chowla_str}
For any $k\in\N$ and any bounded $a\colon\N^k\to\C$, we have
$$
\Esub{n\in[N]} a\big(\Omega(n),\Omega(n+1),\ldots,\Omega(n+k-1)\big)
= \Esub{(n_1,\ldots,n_k)\in[N]^k} a\big(\Omega(n_1),\ldots,\Omega(n_k)\big) + \oh_{N\to\infty}(1).
$$
\end{conjecture}

Conjectures \ref{conj_independence_almost_primes} and \ref{conj_functional_chowla_str} are equivalent. The proof of this equivalence is analogous to the proof that \cref{MT} and \cref{thm_C} are equivalent given in \cref{Proof_Thm_C}, so it is omitted.

\vspace{2mm} 
\noindent
\textbf{Acknowledgments.}
The authors were supported by the Swiss National
Science Foundation grant TMSGI2-211214.
They also thank Harald A. Helfgott for his valuable comments and suggestions on an earlier draft of this article.


%As a final result, we obtain a generalization of the theorem of \Matomaki{}-\Radziwill{}-Tao, which we believe is of independent interest, and it asserts that for any bounded $a\colon\N\to\C$, the function $a(\Omega(n))$ is locally $2$-Fourier uniform.

%\begin{maintheorem}
%\label{thm_local_2-fourier_uniformity}
%For any bounded $a\colon\N\to\C$ we have
%\begin{align}\label{eqn_local_2-fourier_uniformity}
%     & \lim_{H\to\infty}\limsup_{N\to\infty}\sup_{\alpha\in[0,1]}\Esub{n\in[N]}
%     \Bigg|\frac{1}{H}\sum_{h=1}^H a(\Omega(n+h))e(h\alpha)\notag \\
%     & \hspace*{7cm} - 
%     \Bigg(\frac{1}{N}\sum_{m=1}^N a(\Omega(m))\Bigg)
%     \Bigg(\frac{1}{H}\sum_{h=1}^H e(h\alpha)\Bigg)\Bigg|
%     = 0.
%\end{align}
%\end{maintheorem}

%Note that \cref{thm_local_2-fourier_uniformity_lio} corresponds to the case $a(n)=(-1)^n$ of \cref{thm_local_2-fourier_uniformity}.

%To derive \cref{thm_local_2-fourier_uniformity} from \cref{thm_A} one can use a  pretty standard simple argument. The proof is omitted, but we sketch how this argument goes: First we may assume that the averages of $a(\Omega(n))$ are $0$, so that we can get rid of the second term inside the absolute value in \eqref{eqn_local_2-fourier_uniformity}. Then, we note that \eqref{eqn_local_2-fourier_uniformity} is equivalent to the same assertion with the absolute value squared. Expanding the square, rearranging appropriately, and applying \cref{thm_A} for the function $(n,m)\mapsto a(n)\overline{a}(m)$ will conclude the proof.



%SECTION
\section{Preparatory results}\label{prep_lemmas}

In this section we state and prove preliminary results that we will use to prove Theorems~\ref{MT} and~\ref{thm_C}.
The first lemma we need is the following very simple bound that is used multiple times throughout.
\begin{lemma}\label{int_bound}
    For any $\epsilon>0$ and any $C_\epsilon \gg \log\frac{1}{\epsilon}$, we have
    $\displaystyle\int_{C_\epsilon}^\infty e^{-\frac{u^2}{2}} \d u
    \ll \epsilon.$
\end{lemma}

\begin{proof}
    Let $\epsilon>0$, $C_\epsilon = \log\frac{1}{\epsilon}$. Then we have
    $$\int_{C_\epsilon}^\infty e^{-\frac{u^2}{2}} \d u
    \leq \frac{1}{C_\epsilon}\int_{C_\epsilon}^\infty u e^{-\frac{u^2}{2}} \d u
    = \frac{e^{-\frac{C_\epsilon^2}{2}}}{C_\epsilon}
    \ll e^{-C_\epsilon}
    = \epsilon.$$
\end{proof}

\subsection{Passing from \Cesaro{} to logarithmic averages}

The following lemma gives a well-known identity relating \Cesaro{} averages to logarithmic ones.

\begin{lemma}\label{Cesaro_to_log}
Let $f\colon\N\to\C$ be $1$-bounded. Then for any $N\in\N$ and $\epsilon\in\left(0,\frac{1}{2}\right)$, we have
$$\logE_{n\in[N]} f(n) = \logE_{M\in(N^\epsilon,N)}\Esub{n\in[M]} f(n) + \Oh\bigg(\frac{1}{\log{N}} + \epsilon\bigg).$$
\end{lemma}

\begin{proof}
By partial summation, we have 
\begin{align*}
    \sum_{n\leq N} \frac{f(n)}{n} 
    & = \sum_{M<N}\frac{1}{M(M+1)}\sum_{n\leq M} f(n) + \Esub{n\in[N]} f(n) \\
    & = \sum_{M<N}\frac{1}{M}\Esub{n\in[M]} f(n) + \Esub{n\in[N]} f(n) + \Oh(1),
\end{align*}
hence,
$$\logE_{n\in[N]} f(n) = \logE_{M\in[1,N)}\Esub{n\in[M]} f(n) + \Oh\bigg(\frac{1}{\log{N}}\bigg).$$
Moreover, we have
\begin{align*}
    & \Big|\logE_{M\in(N^\epsilon,N)}\Esub{n\in[M]} f(n) - \logE_{M\in[1,N)}\Esub{n\in[M]} f(n)\Big| \\
    & \hspace*{3cm} \leq \bigg|\frac{1}{\log{N}}\sum_{M\in(N^\epsilon,N)}\frac{1}{M}\Esub{n\in[M]} f(n) - \frac{1}{\log{N}}\sum_{M<N}\frac{1}{M}\Esub{n\in[M]} f(n)\bigg| + \Oh(\epsilon) \\
    & \hspace*{3cm} \leq \frac{1}{\log{N}}\sum_{M\leq N^\epsilon}\frac{1}{M} + \Oh(\epsilon) 
    = \Oh(\epsilon).
\end{align*}
Combining the last two estimates, the lemma follows.
\end{proof}


Recall that given $\ell,N\in\N$ we have defined $\overline{\pi}_\ell(N) = \frac{|\P_\ell\cap [N]|}{N}$. We also define
$$\overline{\pi}^{\log}_\ell(N) = \logE_{n\in[N]}\1_{\P_\ell}(n).$$

We note the following immediate identities that will be invoked several times throughout \cref{Proof_Thm_C}: For any function $a\colon\N\to\C$ and any $N\in\N$,
\begin{equation}\label{eqn_expanding_a(Omega(n))_1}
    a(\Omega(n)) = \sum_{\ell=1}^\infty a(\ell)\1_{\P_\ell}(n),
\end{equation}
and consequently,
\begin{equation}\label{eqn_expanding_a(Omega(n))_2}
    \Esub{n\in[N]} a(\Omega(n)) 
    = \sum_{\ell=1}^\infty a(\ell)\overline{\pi}_\ell(N)
    \quad\text{and}\quad
    \logE_{n\in[N]} a(\Omega(n)) 
    = \sum_{\ell=1}^\infty a(\ell)\overline{\pi}^{\log}_\ell(N).
\end{equation}
   
\begin{lemma}\label{pi_k_estimate}
    For any $A>1$ and $N\in\N$ sufficiently large, we have
    \begin{multline*}
        \overline{\pi}_\ell(N) 
        = f(\ell;\mu_N,\sigma_N)\bigg(1 + \Oh\bigg(\frac{1}{\sqrt{\log\log{N}}}\bigg)\bigg) \\
        = \frac{\exp\Big(-\frac{1}{2}\Big(\frac{\ell-\log\log{N}}{\sqrt{\log\log{N}}}\Big)^2\Big)}{\sqrt{2\pi\log\log{N}}}
    \bigg(1 + \Oh\bigg(\frac{1}{\sqrt{\log\log{N}}}\bigg)\bigg),
    \end{multline*}
    uniformly for all $\ell\in T_{A,N}$, where $T_{A,N}$ is the typical range defined in \eqref{typical_range}, and $f(\ell;\mu_N,\sigma_N)$ is the Gaussian probability density function defined in \eqref{eqn_Gaussian_prob_dens_fctn} with mean $\mu_N$ and standard deviation $\sigma_N$ defined in \eqref{mu_N and sigma_N}.
\end{lemma}

\begin{proof}
    This follows immediately from the estimate
    $$\overline{\pi}_\ell(N) = \frac{\exp\Big(-\frac{1}{2}\Big(\frac{\ell-\log\log{N}}{\sqrt{\log\log{N}}}\Big)^2\Big)}{\sqrt{2\pi\log\log{N}}}
    \bigg(1 + \Oh\bigg(\frac{1}{\sqrt{\log\log{N}}} + \frac{|\ell-\log\log{N}|^3}{(\log\log{N})^2}\bigg)\bigg),$$
    which holds uniformly for all positive integers $\ell\leq \frac{3}{2}\log\log{N}$. For this estimate we refer the reader to \cite[p. 236]{MV_book}.
\end{proof}

\begin{lemma}\label{pk_and_logpk}
    For $N\in\N$ sufficiently large, we have  
    $$\sum_{\ell=1}^\infty \big|\overline{\pi}_\ell(N) - \overline{\pi}^{\log}_\ell(N)\big|
    \ll \frac{(\log\log\log{N})^2}{\sqrt{\log\log{N}}}.$$
\end{lemma}


\begin{proof}
    Let $C_N = \log\log\log{N}$ and then by \cref{int_bound}
    we have
    $$\frac{1}{\sqrt{2\pi}}\int_{-C_N}^{C_N} e^{-\frac{u^2}{2}} \d u 
    = 1 + \Oh\bigg(\frac{1}{\sqrt{\log\log{N}}}\bigg).$$
    We define $T_N = (\log\log{N} - C_N\sqrt{\log\log{N}}, \log\log{N} + C_N\sqrt{\log\log{N}}]\cap\Z$ and we denote its complement in $\Z$ by $T_N^\text{c}$. Then, by the quantitative \Erdos-Kac theorem \cite[Theorem 2]{renyi-turan}, we have
    $$\Esub{n\in[N]} \1_{T_N^\text{c}}(\Omega(n))
    = 1 - \frac{1}{\sqrt{2\pi}}\int_{-C_N}^{C_N} e^{-\frac{u^2}{2}} \d u + \Oh\bigg(\frac{1}{\sqrt{\log\log{N}}}\bigg)
    \ll \frac{1}{\sqrt{\log\log{N}}}.$$
    In view of \cref{Cesaro_to_log}, we also have
    $$\logE_{n\in[N]} \1_{T_N^\text{c}}(\Omega(n))
    \ll \frac{1}{\sqrt{\log\log{N}}}.$$
    By the two last estimates, it follows that
    $$\sum_{\ell\in T_N^\text{c}} \overline{\pi}_\ell(N)
    = \Esub{n\in[N]}\sum_{\ell\in T_N^\text{c}} \1_{\P_\ell}(n)
    = \Esub{n\in[N]} \1_{T_N^\text{c}}(\Omega(n))
    \ll \frac{1}{\sqrt{\log\log{N}}}$$
    and respectively,
    $$\sum_{\ell\in T_N^\text{c}} \overline{\pi}^{\log}_\ell(N)
    \ll \frac{1}{\sqrt{\log\log{N}}}.$$
    Therefore, using triangle inequality, we see that it is enough to prove
    \begin{equation*}
        \sum_{\ell\in T_N}\big|\overline{\pi}_\ell(N) - \overline{\pi}^{\log}_\ell(N)\big|
        \ll \frac{(\log\log\log{N})^2}{\sqrt{\log\log{N}}}.
    \end{equation*}
    Applying \cref{Cesaro_to_log} with $\epsilon = (\log\log{N})^{-2}$, we have that
    $$\sum_{\ell\in T_N}\big|\overline{\pi}_\ell(N) - \overline{\pi}^{\log}_\ell(N)\big|
    \leq \sup_{N^\epsilon<M<N}\sum_{\ell\in T_N}\big|\overline{\pi}_\ell(N) - \overline{\pi}_\ell(M)\big| + \Oh\Bigg(\frac{\log\log\log{N}}{(\log\log{N})^\frac{3}{2}}\Bigg),
    $$
    hence, in view of \cref{pi_k_estimate}, it suffices to show that 
    \begin{equation}\label{restriction_TN}
        \sum_{\ell\in T_N}|f(\ell;\mu_N,\sigma_N) - f(\ell;\mu_M,\sigma_M)| \ll \frac{(\log\log\log{N})^2}{\sqrt{\log\log{N}}}
    \end{equation}
    uniformly over all $M\in(N^\epsilon,N)$.
    To show this, we let $M$ be an integer in this range, and we write it as $M = N - K$, for some positive integer $K<N-N^\epsilon$. Now setting $L=\frac{K}{N}<1$, we have
    $$\log\log{M} = \log\log(N-K) = \log(\log{N}+\log(1-L))
    = \log\log{N} + \log\bigg(1 + \frac{\log(1-L)}{\log{N}}\bigg),$$
    where $N^{\epsilon-1}<1-L<1$, hence 
    $$\epsilon-1 = \frac{1}{(\log\log{N})^2}-1 < \frac{\log(1-L)}{\log{N}}<0,$$
    which implies
    $\log\log{M} = \log\log{N} - \eta$, for some $0<\eta<2\log\log\log{N}$.
    Given $t\in\R$, we define the function
    $$f_t(u) = f(t;u;u)
    = \frac{\exp\Big(-\frac{1}{2}\Big(\frac{t-u}{\sqrt{u}}\Big)^2\Big)}{\sqrt{2\pi u}}$$
    and we set $x=\log\log{M}$, $y=\log\log{N}$. Then the left-hand side of \eqref{restriction_TN} becomes
    $$\sum_{\ell\in T_N}|f_\ell(x) - f_\ell(y)|.$$
    By the mean value theorem of calculus, there exists $x_0\in(x,y)$
    such that $|f_\ell(x) - f_\ell(y)| = |f_\ell'(x_0)||x-y| = |f_\ell'(x_0)|\eta$. We have 
    $$f_\ell'(u) = \frac{\ell^2-u^2-u}{2u^2}f_\ell(u),$$
    thus
    \begin{align*}
    |f_\ell(x)-f_\ell(y)| = \frac{|\ell^2-x_0^2-x_0|}{2x_0^2}f_\ell(x_0)\eta
    & \leq \frac{|\ell^2-x_0^2|}{2x_0^2}f_\ell(x_0)\eta + \frac{1}{2x_0}f_\ell(x_0)\eta \\
    & \ll \frac{\log\log\log{N}}{(\log\log{N})^2}|\ell^2-x_0^2|f_\ell(x_0)
    + \frac{\log\log\log{N}}{\log\log{N}}f_\ell(x_0).
    \end{align*}
    If $\ell\geq x_0$, then 
    \begin{multline*}
        \ell^2-x_0^2 \leq (\log\log{N}+C_N\sqrt{\log\log{N}})^2-(\log\log{N}-\eta)^2 \\
        \ll C_N(\log\log{N})^\frac{3}{2} = (\log\log\log{N})(\log\log{N})^\frac{3}{2},
    \end{multline*}
    otherwise, a similar calculation gives the same bound. In any case, we have that
    $$|f_\ell(x)-f_\ell(y)| \ll \frac{(\log\log\log{N})^2}{\sqrt{\log\log{N}}}f_\ell(x_0),$$
    thus, the left-hand side of \eqref{restriction_TN} is 
    $$\ll \frac{(\log\log\log{N})^2}{\sqrt{\log\log{N}}} \sum_{\ell\in T_N}f_\ell(x_0).$$
    Observe that $x_0\sim y$, thus, $f_\ell(x_0)\ll f_\ell(y) = f(\ell;\mu_N,\sigma_N)$ and then
    $$\sum_{\ell\in T_N} f_\ell(x_0) \ll \int_{-\infty}^\infty f(t;\mu_N,\sigma_N) \d t = 1,$$
    which concludes the proof of \eqref{restriction_TN}, and consequently the proof of the lemma.
\end{proof}


\begin{lemma}\label{aOmega_and_logaOmega}
For any bounded function $a\colon \N\to\C$ and $N\in\N$ sufficiently large, we have  
$$
\Big|\Esub{n\in[N]} a(\Omega(n))-\logE_{n\in[N]} a(\Omega(n))\Big| 
\ll_{\|a\|_\infty} \frac{(\log\log\log{N})^2}{\sqrt{\log\log{N}}}.
$$
\end{lemma}

\begin{proof}
Invoking \eqref{eqn_expanding_a(Omega(n))_2}, we have
$$
\Big|\Esub{n\in[N]} a(\Omega(n))-\logE_{n\in[N]} a(\Omega(n))\Big|
=
\Bigg|\sum_{\ell=1}^\infty a(\ell)\big(
\overline{\pi}_\ell(N) -\overline{\pi}^{\log}_\ell(N)\big)
\Bigg|
\ll
\sum_{\ell=1}^\infty \big|
\overline{\pi}_\ell(N) -\overline{\pi}^{\log}_\ell(N)\big|.
$$
The claim now follows from \cref{pk_and_logpk}.
\end{proof}


\subsection{Independence of periodic functions}

We say that a function $f\colon\N\to\C$ can be written as a \emph{linear combination of $k$ indicator functions of infinite arithmetic progressions} if there exist $v_1,\ldots,v_k\in\{z\in\C: |z|\leq 1\}$, $a_1,\ldots,a_k\in\N$ and $b_1,\ldots,b_k\in\N\cup\{0\}$ such that
$$
f(n)=
v_1 \1_{a_1\N+b_1}(n)+\ldots+v_k \1_{a_k\N+b_k}(n),
%\sum_{i=1}^k c_i \1_{a_i\N+b_i}(n),
\qquad\forall n\in\N.
$$
Note that the modulus of the coefficients $v_1,\ldots,v_k$ in the above description is not allowed to exceed $1$.


\begin{lemma}
\label{lem_independence_periodic_fctns_np_1}
Let $k,\ell,r,s\in\N$. Let $f\colon\N\to\C$ be an $r$-periodic function that can be written as a linear combination of $k$ indicator functions of infinite arithmetic progressions and $g\colon \N\to\C$ an $s$-periodic function that can be written as a linear combination of $\ell$ indicator functions of infinite arithmetic progressions. If $r$ and $s$ are coprime, then for any $N\in\N$,
\begin{equation}
\label{eqn_periodic_fctn_independence_np_1}
\Esub{n\in[N]} f(n)g(n) = \Big(\Esub{n\in[N]} f(n)\Big)\Big(\Esub{n\in[N]} g(n)\Big) + \Oh\bigg(\frac{k\ell}{N}\bigg),
\end{equation}
where the implicit constant in $\Oh(k\ell/N)$ is universal (i.e., does not depend on $k$, $\ell$, $r$, $s$, $f$, $g$, or $N$).
\end{lemma}

\begin{proof}
Let  
$f(n)=v_1 \1_{a_1\N+b_1}(n)+\ldots+v_k \1_{a_k\N+b_k}(n)$ and $g(n)=w_1 \1_{c_1\N+d_1}(n)+\ldots+w_\ell \1_{c_\ell\N+d_\ell}(n)$.
Then
\begin{align}
\label{eqn_periodic_fctn_independence_01}
\Esub{n\in[N]} f(n)g(n)
= \sum_{\substack{1\leq i \leq k\\ 1\leq j \leq \ell}} v_i w_j \Big(\Esub{n\in[N]} \1_{a_i\N+b_i}(n)\1_{c_j\N+d_j}(n)\Big).
\end{align}
Since $a_i$ divides $r$ and $c_j$ divides $s$, we get $\gcd(a_i,c_j)=1$. So 
$$\Esub{n\in[N]} \1_{a_i\N+b_i}(n)\1_{c_j\N+d_j}(n)= \frac{1}{a_i c_j} + \Oh\left(\frac{1}{N}\right).$$ 
On the other hand, we have 
$$\Esub{n\in[N]} \1_{a_i\N+b_i}(n) = \frac{1}{a_i}+\Oh\left(\frac{1}{N}\right)
\quad\text{and}\quad
\Esub{n\in[N]} \1_{c_j\N+d_j}(n) = \frac{1}{c_j} + \Oh\left(\frac{1}{N}\right).$$
This implies that
\begin{align}\label{eqn_periodic_fctn_independence_02}
\Esub{n\in[N]} \1_{a_i\N+b_i}(n)\1_{c_j\N+d_j}(n)
= \Big(\Esub{n\in[N]} \1_{a_i\N+b_i}(n)\Big)\Big(\Esub{n\in[N]} \1_{c_j\N+d_j}(n)\Big) +\Oh\left(\frac{1}{N}\right).
\end{align}
Combining \eqref{eqn_periodic_fctn_independence_01} and \eqref{eqn_periodic_fctn_independence_02} gives
\begin{align*}\label{eqn_periodic_fctn_independence_03}
\Esub{n\in[N]} f(n)g(n)
&=
\sum_{\substack{1\leq i \leq k\\ 1\leq j \leq \ell}} v_i w_j \Big(\Esub{n\in[N]} \1_{a_i\N+b_i}(n)\Big)\Big(\Esub{n\in[N]} \1_{c_j\N+d_j}(n)\Big) +\Oh\left(\frac{k\ell}{N}\right)
\\
&=
\Big(\Esub{n\in[N]} f(n)\Big)\Big(\Esub{n\in[N]} g(n)\Big)+\Oh\left(\frac{k\ell}{N}\right),
\end{align*}
as was to be shown. 
\end{proof}


\subsection{\Turan{}-Kubilius inequality and applications}

The following is a special case of the \Turan{}-Kubilius inequality, which we will use to derive rough upper bounds on the distribution of $\Omega(n)$ outside the typical range. Let $n\mapsto \1_{p\mid n}$ denote the indicator function of the set of numbers divisible by a fixed prime $p$.

\begin{theorem}[\Turan{}-Kubilius inequality; see \cite{elliott_TK}]{}
\label{nam_turan_kubilius}
For any $N\in\N$ and $P\subset\P\cap [1,N]$, we have
\begin{equation}
\label{eqn_turan}
\Esub{n\in[N]}\Bigg|\sum_{p\in P}\1_{p\mid n}-\sum_{p\in P}\frac{1}{p}\Bigg| \leq2\Bigg(\sum_{p\in P}\frac{1}{p}\Bigg)^{\frac{1}{2}}.
\end{equation}
\end{theorem}


\begin{corollary}
\label{cor_TK_1}
For $N\in\N$ and $D\geq 1$, let $$\Gamma_{N,D}=\Bigg\{n\leq N: \sum_{p\leq N}\1_{p\mid n}\geq \log\log N+ D\sqrt{\log\log{N}}\Bigg\}$$ and $$\Delta_{N,D} = \Bigg\{n\leq N\colon \sum_{p\leq N}\1_{p\mid n} \leq \log\log N - D\sqrt{\log\log{N}}\Bigg\}.$$
Then for $N$ sufficiently large, we have
$$\frac{|\Gamma_{N,D}|}{N} =\Oh\bigg(\frac{1}{D}\bigg) \quad\text{and}\quad
\frac{|\Delta_{N,D}|}{N} =\Oh\bigg(\frac{1}{D}\bigg).$$
\end{corollary}

\begin{proof}
We prove the result only for $\Gamma_{N,D}$, since the proof for $\Delta_{N,D}$ is analogous.
By \cref{nam_turan_kubilius} and Markov's inequality, we have for all $t>0$ and all but finitely many $N$ that
\[
\frac{1}{N}\Bigg|\Bigg\{n\leq N: \Bigg|\sum_{p\leq N}\1_{p\mid n}-\sum_{p\leq N}\frac{1}{p}\Bigg| >t \Bigg\}\Bigg|
\leq \frac{2}{t} \Bigg(\sum_{p\leq N}\frac{1}{p}\Bigg)^{\frac{1}{2}}
\leq \frac{4\sqrt{\log\log{N}}}{t}.
\]
It follows that
\[
\frac{1}{N}\Bigg|\Bigg\{n\leq N: \sum_{p\leq N}\1_{p\mid n} \geq \sum_{p\leq N}\frac{1}{p} + t \Bigg\}\Bigg|\leq  \frac{4\sqrt{\log\log{N}}}{t}.
\]
Since, by Mertens' theorem, for all but finitely many $N$ we have
\[
\sum_{p\leq N}\frac{1}{p} + \frac{D}{2}\sqrt{\log\log{N}}\leq \log\log N+ D\sqrt{\log\log{N}},
\]
we can take $t=\frac{D}{2}\sqrt{\log\log{N}}$ and get
\[
\frac{|\Gamma_{N,D}|}{N} =\Oh\bigg(\frac{1}{D}\bigg)
\]
as desired.
\end{proof}

\begin{lemma}
\label{lem_approx_deg_K}
Let $N\in\N$ and $M\in\N$ such that $N^\frac{1}{\log\log{N}}\leq M\leq N^{2\log\log{N}}$. Let also $t,\sigma\in\R$ with $\sigma\geq \frac{1}{2}$, $P\subset \P\cap [1,N]$, and $p,q,K\leq H$ be positive integers. Consider the degree $K$ Taylor approximation of $e(x)$, given by
\[
e_K(x)=\sum_{k=0}^K \frac{(2\pi i)^k x^k}{k!}.
\]
If $K\geq 200 t (\log\log N)^\sigma$, then for $N$ sufficiently large, we have
\begin{multline}
\label{eqn_approx_deg_K_1}
\Esub{n\in[N]} e\bigg(\frac{t\left(\sum_{r\in P}\1_{r\mid (n+p)}-\1_{r\mid (n+q)}\right)}{\sqrt{\log\log{N}}}\bigg)
\\
=\Esub{n\in[N]}e_K\bigg(\frac{t\left(\sum_{r\in P}\1_{r\mid (n+p)}-\1_{r\mid (n+q)}\right)}{\sqrt{\log\log{N}}}\bigg) + \Oh\bigg(e^{-K}+\frac{1}{(\log\log N)^\sigma}+\frac{H}{N}\bigg)
\end{multline}
and
\begin{multline}
\label{eqn_approx_deg_K_2}
\Esub{\substack{m\in[M] \\ n\in[N]}}e\bigg(\frac{t\left(\sum_{r\in P}\1_{r\mid m}-\1_{r\mid n}\right)}{\sqrt{\log\log{N}}}\bigg)
\\
=\Esub{\substack{m\in[M] \\ n\in[N]}}e_K\bigg(\frac{t\left(\sum_{r\in P}\1_{r\mid m}-\1_{r\mid n}\right)}{\sqrt{\log\log{N}}}\bigg) + \Oh\bigg(e^{-K}+\frac{1}{(\log\log N)^\sigma}\bigg).
\end{multline}
\end{lemma}

\begin{proof}
We only provide a proof of \eqref{eqn_approx_deg_K_1}. The proof of \eqref{eqn_approx_deg_K_2} is analogous. Note that
\begin{align*}
|e(x)-e_K(x)|
=\bigg|\sum_{k=K+1}^\infty \frac{(2\pi i)^k x^k}{k!}\bigg|
& = \bigg|(2\pi i)^{K+1} x^{K+1}\sum_{k=0}^\infty \frac{(2\pi i)^{k} x^{k}}{(k+K+1)!}\bigg|
\\
&\leq \bigg|\frac{(2\pi i)^{K+1} x^{K+1}}{(K+1)!}\bigg| \bigg|\underbrace{\sum_{k=0}^\infty \frac{(2\pi i)^{k} x^{k}}{k!}}_{=1}\bigg|.
\end{align*}
Let $D=(\log\log N)^\sigma$, and let $\Gamma_{N,D}$ be as in \cref{cor_TK_1}. Then, by the conclusion of \cref{cor_TK_1}, we have
\begin{align*}
& \Esub{n\in[N]} e\bigg(\frac{t\left(\sum_{r\in P}\1_{r\mid (n+p)}-\1_{r\mid (n+q)}\right)}{\sqrt{\log\log{N}}}\bigg) 
\\
& \hspace*{.5cm} = \Esub{n\in[N]} \1 _{\Gamma_{N,D}}(n+p)\1 _{\Gamma_{N,D}}(n+q) e\bigg(\frac{t\left(\sum_{r\in P}\1_{r\mid (n+p)}-\1_{r\mid (n+q)}\right)}{\sqrt{\log\log{N}}}\bigg) \\
& \hspace*{4cm} + \Oh\bigg(\frac{1}{(\log\log N)^\sigma}+\frac{H}{N}\bigg)
\\
& \hspace*{.5cm} = \Esub{n\in[N]} \1 _{\Gamma_{N,D}}(n+p)\1 _{\Gamma_{N,D}}(n+q) e_K\bigg(\frac{t\left(\sum_{r\in P}\1_{r\mid (n+p)}-\1_{r\mid (n+q)}\right)}{\sqrt{\log\log{N}}}\bigg)
\\
& \hspace*{4cm} + \Oh\bigg(\sup_{n\in \Gamma_{N,D}} \frac{20^{K+1} t^{K+1} \left(\sum_{r\in P}\1_{r\mid n}\right)^{K+1} }{(K+1)!(\log\log N)^{\frac{K+1}{2}}} +\frac{1}{(\log\log N)^\sigma}+\frac{H}{N}\bigg)
\\
& \hspace*{.5cm} = \Esub{n\in[N]}  e_K\bigg(\frac{t\left(\sum_{r\in P}\1_{r\mid (n+p)}-\1_{r\mid (n+q)}\right)}{\sqrt{\log\log{N}}}\bigg)
\\
&\hspace*{4cm} + \Oh\bigg(\sup_{n\in \Gamma_{N,D}} \frac{20^{K+1} t^{K+1} \left(\sum_{r\in P}\1_{r\mid n}\right)^{K+1} }{(K+1)!(\log\log N)^{\frac{K+1}{2}}} +\frac{1}{(\log\log N)^\sigma}+\frac{H}{N}\bigg).
\end{align*}
We now have
\begin{align*}
\sup_{n\in \Gamma_{N,D}} \frac{20^{K+1} t^{K+1} \left(\sum_{r\in P}\1_{r\mid n}\right)^{K+1} }{(K+1)!(\log\log N)^{\frac{K+1}{2}}} 
&\leq 
\frac{40^{K+1} t^{K+1} (\log\log N)^{(K+1)\sigma} }{(K+1)!}.  
\end{align*}
Note, by Stirling's formula, we have $\log(m!)\geq m\log m-2m$ for all but finitely many $m$, and hence for all but finitely many $N$ we have that
\begin{align*}
&\log\bigg(\frac{40^{K+1} t^{K+1} (\log\log N)^{(K+1)\sigma} }{(K+1)!}\bigg)
\\
&\qquad\qquad\leq (K+1)\log(40t)+ (K+1)\sigma(\log\log\log N)- (K+1)\log(K)+2(K+1)
\\
&\qquad\qquad\leq (K+1)\left(\log(40t)+ \sigma (\log\log\log N)- \log(K)+2\right)
\\
&\qquad\qquad= (K+1)\bigg(\log\bigg(\frac{40 t (\log\log N)^\sigma}{K}\bigg)+2\bigg).
\end{align*}
Invoking the assumption $K\geq 200 t (\log\log N)^\sigma$, it now follows that
\[
\log\bigg(\frac{40^{K+1} t^{K+1} (\log\log N)^{\sigma(K+1)} }{(K+1)!}\bigg)\leq -(K+1)
\]
and hence
\begin{align*}
\sup_{n\in \Gamma_{N,D}} \frac{20^{K+1} t^{K+1} \left(\sum_{r\in P}\1_{r\mid n}\right)^{K+1} }{(K+1)!(\log\log N)^{\frac{K+1}{2}}} 
&\leq e^{-(K+1)}.
\end{align*}
The conclusion now follows by combining the above estimates.
\end{proof}

\subsection{Approximations of $\Omega(n)$}

For technical reasons, it will be convenient to replace $\Omega(n)$ with an approximation $\omega_N(n)$ that ignores the contribution of large primes.
We define 
\begin{equation}\label{omega_N_def}
     t_N=N^{\frac{1}{(\log\log N)^8}}
     \quad\text{and}\quad
     \omega_N(n)=\sum_{\substack{p\in\P \\ p\leq t_N}} \1_{p\mid n}.
\end{equation}

The following lemma estimates the error that occurs when replacing $\Omega(n)$ with $\omega_N(n)$ in the expressions appearing in the proof of our main theorem.

\begin{lemma}
\label{lem_approx_Omega}
Let $N,H\in\N$ and $p,q\leq H$ be positive integers.
For any $t\in\R$, we have
\begin{multline}
\label{eqn_approx_Omega_1}
\Esub{n\in[N]} e\bigg(\frac{t\left(\Omega(n+p)-\Omega(n+q)\right)}{\sqrt{\log\log{N}}} \bigg)
\\
=\Esub{n\in[N]}  e\bigg(  \frac{t\left(\omega_N(n+p)-\omega_N(n+q)\right)}{\sqrt{\log\log{N}}} \bigg)+\Oh\left(\frac{t(\log\log\log N)}{\sqrt{\log\log{N}}}+\frac{tH}{N}\right)
\end{multline}
and
\begin{equation}
\label{eqn_approx_Omega_2}
\Esub{\substack{m\in[M] \\ n\in[N]}} e\bigg(  \frac{t\left(\Omega(m)-\Omega(n)\right)}{\sqrt{\log\log{N}}} \bigg)=\Esub{\substack{m\in[M] \\ n\in[N]}}  e\bigg(  \frac{t\left(\omega_N(m)-\omega_N(n)\right)}{\sqrt{\log\log{N}}} \bigg)+\Oh\left(\frac{t(\log\log\log N)}{\sqrt{\log\log{N}}}\right).
\end{equation}
\end{lemma}

\begin{proof}
We only prove \eqref{eqn_approx_Omega_1}, as the proof of \eqref{eqn_approx_Omega_2} is analogous.
Using that $|e(x)-1|^2 = 2-2\cos{x} \leq x^2$ for all $x\in\R$, we get
\begin{align*}
&\Esub{n\in[N]} \left|e\bigg(  \frac{t\left(\Omega(n+p)-\Omega(n+q)\right)}{\sqrt{\log\log{N}}} \bigg)-e\bigg(  \frac{t\left(\omega_N(n+p)-\omega_N(n+q)\right)}{\sqrt{\log\log{N}}} \bigg)\right|
\\
&\hspace*{1.3cm} = \Esub{n\in[N]} \left|e\bigg( \frac{t \left|\left(\Omega(n+p)-\Omega(n+q)\right) - \left(\omega_N(n+p)-\omega_N(n+q)\right)\right|}{\sqrt{\log\log{N}}}\bigg)-1\right|
\\
&\hspace*{1.3cm} \leq \frac{t}{\sqrt{\log\log{N}}}
\Esub{n\in[N]}\left|\left(\Omega(n+p)-\Omega(n+q)\right) - \left(\omega_N(n+p)-\omega_N(n+q)\right)\right|
\\
&\hspace*{1.3cm}
\ll \frac{t}{\sqrt{\log\log{N}}}\Esub{n\in[N]}\left|\Omega(n)-\omega_N(n)\right| + \frac{tH}{N}.
\end{align*}
Now, we have
\begin{multline*}
\Esub{n\in[N]}\left|\Omega(n)-\omega_N(n)\right|
=\Esub{n\in[N]}\sum_{t_N< p \leq N}\1_{p\mid n}+\Oh(1)
=\sum_{t_N< p \leq N} \frac{1}{N}\left\lfloor \frac{N}{p}\right\rfloor+\Oh(1)
\\
\leq \sum_{t_N< p \leq N} \frac{1}{p} +\Oh(1)
\ll \log\bigg(\frac{\log{N}}{\log{t_N}}\bigg)
= 8\log\log\log{N}.
\end{multline*}
\end{proof}



%SECTION
\section{Reduction to multiplicative functions}\label{reduction_section}

The next three sections are dedicated to proving \cref{MT}. In this section, using a key proposition (\cref{Helf-Radz_cor}), which follows from \cite[Corollary 1.4]{helf_radz}, and a Fourier analytic argument, we reduce the statement in \cref{MT} to an assertion concerning averages that involve a sequence of multiplicative functions $(\f)_{\xi\in I_N}$, where $I_N$ is a set of frequencies (see \cref{reducing_MT}). Then we state two theorems in which we handle separately small (see \cref{small_xi_thm}) and large frequencies (see \cref{large_xi_thm}), which are proved in Sections~\ref{small_xi_section} and \ref{large_xi_section} respectively. 

First, note that in light of \cref{aOmega_and_logaOmega}, we may substitute the \Cesaro{} averages in the left-hand side of \eqref{MT_eq} by logarithmic ones, hence \cref{MT} reduces to establishing that for any bounded functions $a,b\colon\N\to\C$, we have
\begin{equation}\label{MT_eq_with_log}
    \logE_{n\in[N]} a(\Omega(n))b(\Omega(n+1)) 
    = \Big(\logE_{n\in[N]}a(\Omega(n))\Big)\Big(\logE_{n\in[N]}b(\Omega(n)\Big)
    + \Oh\Bigg(\frac{\sqrt{\log\log\log{N}}}{(\log\log{N})^\frac{1}{6}}\Bigg)
\end{equation}
Note that by normalizing, it is enough to prove the above for $1$-bounded functions. Throughout the next three sections we fix two $1$-bounded functions $a,b\colon\N\to\C$ and $N\in\N$ sufficiently large (depending on any absolute constant that will arise).
We define 
\begin{equation}\label{I_N_def}
    A = A(N) = 4(\log\log{N})^{\frac{1}{9}}
    \quad\text{and}\quad
    I_N = (-A\sqrt{\log\log{N}},A\sqrt{\log\log{N}}]\cap\Z.
\end{equation}

The following lemma is a consequence of R\'enyi-Tur\'an's quantitative \Erdos-Kac theorem (see \cite[Theorem 2]{renyi-turan}).

\begin{lemma}\label{renyi-turan_cor}
    We have
    $$\frac{|\{n\leq N\colon \Omega(n)>\log\log{N}+A\sqrt{\log\log{N}}\}|}{N}
    \ll \frac{1}{{\sqrt{\log\log{N}}}}$$
    and
    $$\frac{|\{n\leq N\colon \Omega(n)<\log\log{N}-(A-1)\sqrt{\log\log{N}}\}|}{N}
    \ll \frac{1}{{\sqrt{\log\log{N}}}}.$$    
\end{lemma}

\begin{proof}
We prove the first bound, as the proof of the other is analogous.
By \cite[Theorem 2]{renyi-turan}, we have
$$\frac{|\{n\leq N\colon \Omega(n)>\log\log{N}+A\sqrt{\log\log{N}}\}|}{N}
= \frac{1}{\sqrt{2\pi}}\int_A^\infty e^{-\frac{u^2}{2}}\d u + \Oh\bigg(\frac{1}{\sqrt{\log\log{N}}}\bigg).$$
The claim follows by invoking \cref{int_bound}.
\end{proof}

For $\xi\in I_N$, we define the completely multiplicative function 
\begin{equation}\label{defn_f}
    \f(n) = e\bigg(\frac{\xi\Omega(n)}{|I_N|}\bigg)
\end{equation}
and for simplicity, we define
$$\Omega_N(n) = \frac{\Omega(n)-\log\log{N}}{\sqrt{\log\log{N}}}.$$
We let
\begin{equation}\label{defn_H0}
    H_0 = H_0(N) = \exp\Big(\exp\big(-(\log\log{N})^{\frac{1}{3}}\big)(\log{N})^\frac{1}{(\log\log{N})^{4/9}}\Big),
\end{equation}
\begin{equation}\label{defn_H}
    H = H(N) = H_0^{\exp((\log\log{N})^{1/3})}
    = \exp\Big((\log{N})^\frac{1}{(\log\log{N})^{4/9}}\Big),
\end{equation}
and finally,
\begin{equation}\label{defn_P}
    \Pc_N = \P\cap[H_0,H].
\end{equation}
Note that
\begin{equation}\label{eqn_L}
    \mathscr{L} = \mathscr{L}_N := \sum_{p\in\Pc_N}\frac{1}{p} \sim (\log\log{N})^\frac{1}{3}.
\end{equation}

The next key proposition follows from \cite[Corollary 1.4]{helf_radz}.
\begin{proposition}\label{Helf-Radz_cor}
    We have
    \begin{equation*}
        \logE_{n\in[N]} a(\Omega(n))b(\Omega(n+1))
        = \logE_{p\in\Pc_N}\logE_{n\in[N]} a(\Omega(n)-1)b(\Omega(n+p)-1) + \Oh\Bigg(\frac{1}{(\log\log{N})^\frac{1}{6}}\Bigg).
    \end{equation*}
\end{proposition}

\begin{proof}[Proof of \cref{Helf-Radz_cor}]
    First, we have
    \begin{align}\label{eqn_introducing_primes}
        \logE_{n\in[N]} a(\Omega(n)) & b(\Omega(n+1)) \notag \\
        & = \logE_{p\in\Pc_N}\logE_{n\in[N]}  a(\Omega(pn)-1)b(\Omega(pn+p)-1) \notag \\
        & =\logE_{p\in\Pc_N}\logE_{n\in[N]} p\1_{p\mid n}a(\Omega(n)-1)b(\Omega(n+p)-1) + \Oh\bigg(\frac{\log{H}}{\log{N}}\bigg).
    \end{align}
    Observe that for any $\epsilon>0$, we can restrict the logarithmic averages from over $n$ from $[N]$ to $(N^\epsilon,N]$, gaining an extra $\Oh(\epsilon)$ error term. We choose $\epsilon = (\log\log{N})^{-1}$. Now we split the interval $(N^\epsilon,N]$ into dyadic intervals as follows: we let $J_1 = \left\lfloor\frac{\epsilon(\log{N})}{\log{2}}\right\rfloor+1$, $J_2 = \left\lceil\frac{\log{N}}{\log{2}}\right\rceil$ and then 
    $(N^\epsilon,N] = \bigcup_{J_1\leq j \leq J_2}(2^{j-1},2^j]$. We note that uniformly for all $1$-bounded functions $f$, we have
    \begin{align*}
        \Esub{j\in[J_1,J_2]}\logE_{n\in(2^{j-1},2^j]} f(n) 
        & \sim
        \frac{\log{2}}{(1-\epsilon)(\log{N})}\sum_{J_1\leq j\leq J_2}\frac{1}{\log{2}}\sum_{2^{j-1}<n\leq 2^j} f(n) \\
        & = \frac{1}{(1-\epsilon)(\log{N})}\sum_{N^\epsilon<n\leq N} f(n)
        \sim \logE_{n\in(N^\epsilon,N]} f(n),
    \end{align*}
    hence applying this for the sequence of functions $(f_p)_{p\in\Pc_N}$, $f_p(n)=\logE_{n\in(N^\epsilon,N]}p\1_{p\mid n} a(\Omega(n)-1)b(\Omega(n+p)-1)$, then in view of \eqref{eqn_introducing_primes}, we obtain
    \begin{multline}\label{split_into_dyadic}
        \logE_{n\in[N]}a(\Omega(n))b(\Omega(n+1)) \\
        = \Esub{j\in[J_1,J_2]}\logE_{p\in\Pc_N}\logE_{n\in(2^{j-1},2^j]}p\1_{p\mid n} a(\Omega(n)-1)b(\Omega(n+p)-1) + \Oh\bigg(\frac{1}{\log\log{N}}\bigg).
    \end{multline}
   Now we want to remove the factor $p\1_{p\mid n}$ in the main term of \eqref{eqn_introducing_primes}. We will use \cite[Corollary 1.4]{helf_radz} for each $j\in[J_1,J_2]$. So we fix $j$ and we check that the assumptions are satisfied:
    \begin{itemize}
        \item $\log{H_0}\geq(\log{H})^\frac{2}{3}(\log\log{H})^2$:
        For $N$ sufficiently large, we have
        $$\frac{(\log{H})^\frac{2}{3}(\log\log{H})^2}{\log{H_0}}
        \leq \frac{(\log{H})^\frac{3}{4}}{\log{H_0}}
        = \frac{\exp\big((\log\log{N})^\frac{1}{3}\big)}{(\log{N})^\frac{1}{4(\log\log{N})^{4/9}}}\leq 1.$$
        \item $\log{H}\leq \sqrt{(\log{2^{j-1}})/\mathscr{L}}$:
        For $N$ sufficiently large, we have
        $$\frac{(\log{H})^2\mathscr{L}}{\log{2^{j-1}}} 
        \ll \frac{(\log{N})^\frac{2}{(\log\log{N})^{4/9}}(\log\log{N})^\frac{1}{3}}{\epsilon(\log{N})}
        = (\log{N})^{\frac{2}{(\log\log{N})^{4/9}}-1}(\log\log{N})^\frac{4}{3}
        \leq 1.$$
    \end{itemize}
    Then, by \cite[Corollary 1.4]{helf_radz}, it follows that 
    \begin{align*}
        \logE_{p\in\Pc_N} & \logE_{n\in(2^{j-1},2^j]}p\1_{p\mid n}a(\Omega(n)-1)b(\Omega(n+p)-1) \\
        & = \logE_{p\in\Pc_N}\logE_{n\in(2^{j-1},2^j]}a(\Omega(n)-1)b(\Omega(n+p)-1) + \Oh\Bigg(\frac{1}{(\log\log{2^{j-1}})^\frac{1}{6}}\Bigg) \\
        & = \logE_{p\in\Pc_N}\logE_{n\in(2^{j-1},2^j]}a(\Omega(n)-1)b(\Omega(n+p)-1) + \Oh\Bigg(\frac{1}{(\log\log{N})^\frac{1}{6}}\Bigg).
    \end{align*}
    By averaging this over all $j\in[J_1,J_2]$ and combining with \eqref{split_into_dyadic}, the result follows.    
\end{proof}

\begin{theorem}\label{reducing_MT}
We have
\begin{multline*}
    \Big|\logE_{n\in[N]}a(\Omega(n))b(\Omega(n+1)) - \Big(\logE_{n\in[N]}a(\Omega(n))\Big)\Big(\logE_{n\in[N]}b(\Omega(n))\Big)\Big| \\
    \leq \Bigg(\sum_{\xi\in I_N}\logE_{n\in[N]}\Big|\logE_{p\in\Pc_N} \f(n+p) - \logE_{m\in[N]}\f(m)\Big|^2\Bigg)^\frac{1}{2} + \Oh\Bigg(\frac{1}{(\log\log{N})^\frac{1}{6}}\Bigg).
\end{multline*}
\end{theorem}

In view of \cref{reducing_MT}, the next theorem immediately implies \eqref{MT_eq_with_log}, thus, \cref{MT}.

\begin{theorem}\label{MT_reduced}
We have
\begin{equation}\label{MT_reduced_eq}
    \sum_{\xi\in I_N}\logE_{n\in[N]}\Big|\logE_{p\in\Pc_N} \f(n+p) - \logE_{m\in[N]}\f(m)\Big|^2
    \ll
    \frac{\log\log\log{N}}{(\log\log{N})^\frac{1}{3}}.
\end{equation}
\end{theorem}

\cref{reducing_MT} is proved in this section. As for \cref{MT_reduced}, we have to treat differently the case of small and the case of large frequencies $\xi\in I_N$. In particular, we reduce the proof of \cref{MT_reduced} to showing the next two theorems.

\begin{theorem}\label{small_xi_thm}
We have
\begin{equation*}
    \sum_{|\xi|\leq A^{5/2}}
    \logE_{n\in[N]}\Big|\logE_{p\in\Pc_N}\f(n+p) - \logE_{m\in[N]} \f(m)\Big|^2 
    \ll
    \frac{\log\log\log{N}}{(\log\log{N})^\frac{1}{3}}.
\end{equation*}
\end{theorem}

\begin{theorem}\label{large_xi_thm}
We have
\begin{equation*}\label{large_xi_eq}
    \sum_{\substack{\xi\in I_N \\ |\xi|>A^{5/2}}}
    \logE_{n\in[N]}\Big|\logE_{p\in\Pc_N}\f(n+p) - \logE_{m\in[N]} \f(m)\Big|^2 
    \ll \frac{1}{(\log\log{N})^\frac{1}{3}}.
\end{equation*}
\end{theorem}

Theorems~\ref{small_xi_thm} and~\ref{large_xi_thm} are proved in Sections~\ref{small_xi_section} and~\ref{large_xi_section} respectively.

We remark that by choosing an even larger constant $C$ in \eqref{I_N_def}, we can make the error term in \cref{large_xi_thm} much smaller, in fact smaller than $(\log\log{N})^{-D}$ for any positive arbitrarily large constant $D$. This can be readily checked in the proof of \cref{large_xi_thm}. However, the above bound is more than enough for our purposes.

\begin{proof}[Proof of \cref{reducing_MT}]
We start by defining the following functions:
\begin{itemize}
    \item $\widetilde{a},\widetilde{b}\colon\N\to\C$, given by 
    $$\widetilde{a}(n) 
    = 
    \begin{cases}
        a(n-1), & \text{if}~n\geq 2, \\
        0, & \text{otherwise},
    \end{cases}
    \quad\text{and}\quad
    \widetilde{b}(n) 
    = 
    \begin{cases}
        b(n-1), & \text{if}~n,\geq 2, \\
        0, & \text{otherwise}.
    \end{cases}
    $$
    \item $b_N\colon\N\to\C$ is given by
    $$b_N(m) 
    := 
    \begin{cases}
        b(m+\lfloor\log\log{N}\rfloor), & \text{if}~m\in I_N, \\
        0, & \text{otherwise},
    \end{cases}
    $$
    \item $\widetilde{b}_N\colon\N\to\C$ given by
    $$
    \widetilde{b}_N(m) = 
    \begin{cases}
        b_N(m-1), & \text{if}~m\geq 2 \\
        0, & \text{otherwise}
    \end{cases}
    \,=
    \begin{cases}
        \widetilde{b}(m+\lfloor\log\log{N}\rfloor), & \text{if}~m\in I_N+1 \\
        0, & \text{otherwise.}
    \end{cases}
    $$
\end{itemize}
By \cref{Helf-Radz_cor}, we have
\begin{align}\label{applying_Tao}
    & \Big|\logE_{n\in[N]} a(\Omega(n))b(\Omega(n+1)) - \Big(\logE_{n\in[N]}a(\Omega(n))\Big)\Big(\logE_{n\in[N]}b(\Omega(n))\Big)\Big| \notag \\
    & \hspace*{.8cm} = \Big|\logE_{p\in\Pc_N}\logE_{n\in[N]}\widetilde{a}(\Omega(n))\widetilde{b}(\Omega(n+p)) - \Big(\logE_{n\in[N]} \widetilde{a}(\Omega(n))\Big)\Big(\logE_{n\in[N]}\widetilde{b}(\Omega(n))\Big)\Big| \notag \\
    & \hspace*{10cm} + \Oh\Bigg(\frac{1}{(\log\log{N})^\frac{1}{6}}\Bigg).
\end{align}
We now define the interval 
$$J_N = \left(-A + \frac{1-\{\log\log{N}\}}{\sqrt{\log\log{N}}}, A + \frac{1-\{\log\log{N}\}}{\sqrt{\log\log{N}}}\right].$$
By definition, if $\Omega(n)-\lfloor\log\log{N}\rfloor\in I_N+1$,
which is equivalent to $\Omega_N(n)\in J_N$, then
$$\widetilde{b}(\Omega(n)) = \widetilde{b}_N(\Omega(n)-\lfloor\log\log{N}\rfloor),$$
and when $\Omega_N(n)\not\in J_N$, we have $\widetilde{b}_N(\Omega(n)-\lfloor\log\log{N}\rfloor) = 0$, 
thus,
\begin{align}\label{a_to_b}
     &\Big|\logE_{p\in\Pc_N}\logE_{n\in[N]} \widetilde{a}(\Omega(n))\widetilde{b}(\Omega(n+p)) - \Big(\logE_{n\in[N]}\widetilde{a}(\Omega(n))\Big)\Big(\logE_{n\in[N]}\widetilde{b}(\Omega(n))\Big)\Big| \notag \\
     & = \Big|\logE_{p\in\Pc_N}\logE_{n\in[N]}\1_{J_N}(\Omega_N(n+p))\widetilde{a}(\Omega(n))\widetilde{b}_N(\Omega(n+p)-\lfloor\log\log{N}\rfloor) \notag \\
     & \hspace*{1.8cm} - \Big(\logE_{n\in[N]}\widetilde{a}(\Omega(n))\Big)\Big(\logE_{n\in[N]}\1_{J_N}(\Omega_N(n))\widetilde{b}_N(\Omega(n)-\lfloor\log\log{N}\rfloor)\Big)\Big| 
     + \Oh(R_1 + R_2) \notag \\
     & = \Big|\logE_{p\in\Pc_N}\logE_{n\in[N]}\widetilde{a}(\Omega(n))\widetilde{b}_N(\Omega(n+p)-\lfloor\log\log{N}\rfloor) \notag \\
     & \hspace*{1.8cm} - \Big(\logE_{n\in[N]}\widetilde{a}(\Omega(n))\Big)\Big(\logE_{n\in[N]}\widetilde{b}_N(\Omega(n)-\lfloor\log\log{N}\rfloor)\Big)\Big| 
     + \Oh(R_1 + R_2) \notag \\
     & = \Big|\logE_{n\in[N]}\widetilde{a}(\Omega(n))\Big(\logE_{p\in\Pc_N}\widetilde{b}_N(\Omega(n+p)-\lfloor\log\log{N}\rfloor) \notag \\
     & \hspace*{1.8cm} - \logE_{m\in[N]}\widetilde{b}_N(\Omega(m)-\lfloor\log\log{N}\rfloor)\Big)\Big| + \Oh(R_1+R_2),
\end{align}
where 
$$R_1 := \logE_{p\in\Pc_N}\logE_{n\in[N]} \left(\1_{(A,\infty)}\left(\Omega_N(n+p)\right) + \1_{(-\infty,-A+1)}\left(\Omega_N(n+p)\right)\right)$$
and 
$$R_2 := \logE_{m\in[N]} \left(\1_{(A,\infty)}\left(\Omega_N(m)\right) + \1_{(-\infty,-A+1)}\left(\Omega_N(m)\right)\right).$$
By \cref{renyi-turan_cor}, we have 
$$\Esub{n\in[N]} \1_{(A,\infty)}\left(\Omega_N(n)\right) \ll \frac{1}{\sqrt{\log\log{N}}}.$$
By \cref{Cesaro_to_log} applied for $\epsilon=(\log{N})^{-1/2}$ (so that $N^\epsilon\to\infty$), the above gives
\begin{equation}\label{indicator_omega_eq1}
    \logE_{n\in[N]} \1_{(A,\infty)}\left(\Omega_N(n)\right) 
    \ll \frac{1}{\sqrt{\log\log{N}}}.
\end{equation}
For any prime $p\leq H$, using shift-invariance we have
\begin{equation}\label{indicator_omega_eq2}
    \logE_{n\in[N]}\1_{(A,\infty)}\left(\Omega_N(n+p)\right)
    \ll \frac{1}{\sqrt{\log\log{N}}} + \frac{\log{H}}{\log{N}}
    \ll \frac{1}{\sqrt{\log\log{N}}}.
\end{equation}
Similarly, using \cref{renyi-turan_cor} and \cref{Cesaro_to_log} again, we have 
\begin{equation}\label{indicator_omega_eq3}
    \logE_{n\in[N]} \1_{(-\infty,-A+1)}\left(\Omega_N(n)\right) 
    \ll \frac{1}{\sqrt{\log\log{N}}}.
\end{equation}
and for any prime $p\leq H$,
\begin{equation}\label{indicator_omega_eq4}
\logE_{n\in[N]}\1_{(-\infty,-A+1)}\left(\Omega_N(n+p)\right) 
\ll \frac{1}{\sqrt{\log\log{N}}}.
\end{equation}
By \eqref{indicator_omega_eq2} and \eqref{indicator_omega_eq4}, it follows that $R_1\ll\frac{1}{\sqrt{\log\log{N}}}$, 
and by \eqref{indicator_omega_eq1} and \eqref{indicator_omega_eq3}, the same bound applies for $R_2$.
Thus, \eqref{a_to_b} becomes
\begin{align}\label{a_to_b_final}
    & \Big|\logE_{p\in\Pc_N}\logE_{n\in[N]}\widetilde{a}(\Omega(n))\widetilde{b}(\Omega(n+p)) - \Big(\logE_{n\in[N]}\widetilde{a}(\Omega(n))\Big)\Big(\logE_{n\in[N]}\widetilde{b}(\Omega(m))\Big)\Big| \notag \\
    & = \Big|\logE_{p\in\Pc_N}\logE_{n\in[N]}\widetilde{b}_N(\Omega(n+p)-\lfloor\log\log{N}\rfloor) - \logE_{n,m\in[N]}\widetilde{b}_N(\Omega(m)-\lfloor\log\log{N}\rfloor)\Big| \notag \\
    & \hspace*{10.5cm}  
    + \Oh\bigg(\frac{1}{\sqrt{\log\log{N}}}\bigg).
\end{align}
Now we apply the Fourier transform to the function $\widetilde{b}_N$ to find
\begin{align*}
    &\widetilde{b}_N(\Omega(m)-\lfloor\log\log{N}\rfloor)
    = b_N(\Omega(m)-\lfloor\log\log{N}\rfloor-1) \\
    & =\sum_{\xi\in I_N}\widehat{b}_N(\xi)e\bigg(\frac{\xi\left(\Omega(m)-\lfloor\log\log{N}\rfloor-1\right)}{|I_N|}\bigg) 
    = \sum_{\xi\in I_N}\widehat{b}_N(\xi)e\bigg(\frac{\xi\left(\lfloor\log\log{N}\rfloor-1\right)}{|I_N|}\bigg)\f(m).
\end{align*}
Substituting the above in the right-hand side of \eqref{a_to_b_final} and rearranging we obtain

\begin{align}\label{after_fourier_exp_eq}
    & \Big|\logE_{p\in\Pc_N}\logE_{n\in[N]}a(\Omega(n))b(\Omega(n+p)) - \Big(\logE_{n\in[N]}a(\Omega(n))\Big)\Big(\logE_{n\in[N]}b(\Omega(m))\Big)\Big| \notag \\
    & = \Bigg|\sum_{\xi\in I_N} e\bigg(\frac{\xi\left(\lfloor\log\log{N}\rfloor-1\right)}{|I_N|}\bigg)\widehat{b}_N(\xi) \notag \\
    & \hspace*{1.4cm} \logE_{n\in[N]}\widetilde{a}(\Omega(n))
    \Big(\logE_{p\in\Pc_N}\f(n+p) - \logE_{m\in[N]}\f(m)\Big)\Bigg| + \Oh\bigg(\frac{1}{\sqrt{\log\log{N}}}\bigg).
\end{align}
Using Cauchy-Schwarz inequality, the main term of \eqref{after_fourier_exp_eq} is
\begin{align*}
    \Bigg|\sum_{\xi\in I_N} e\bigg(\frac{\xi\left(\lfloor\log\log{N}\rfloor-1\right)}{|I_N|}\bigg)\widehat{b}_N(\xi)\logE_{n\in[N]}\widetilde{a}(\Omega(n))
    \Big(\logE_{p\in\Pc_N}\f(n+p) - \logE_{m\in[N]}\f(m)\Big)\Bigg| & \\
    \leq \Bigg(\sum_{\xi\in I_N}|\widehat{b}_N(\xi)|^2\Bigg)^\frac{1}{2}
    \Bigg(\sum_{\xi\in I_N}\Big|\logE_{n\in[N]}\widetilde{a}(\Omega(n))
    \Big(\logE_{p\in\Pc_N}\f(n+p) - \logE_{m\in[N]}\f(m)\Big)\Big|^2\Bigg)^\frac{1}{2}. &
\end{align*}
By Parseval's identity, we have
$\sum_{\xi\in I_N}|\widehat{b}_N(\xi)|^2
= \Esub{n\in I_N}|b_N(n)|^2
\leq \|b\|_\infty
\leq 1$,
so, using Cauchy-Schwarz inequality again, it follows that
\begin{align}\label{bounding_main_term_eq1}
    & \Bigg|\sum_{\xi\in I_N} e\bigg(\frac{\xi\left(\lfloor\log\log{N}\rfloor-1\right)}{|I_N|}\bigg)\widehat{b}_N(\xi)\logE_{n\in[N]}\widetilde{a}(\Omega(n))
    \Big(\logE_{p\in\Pc_N}\f(n+p) - \logE_{m\in[N]}\f(m)\Big)\Bigg| \notag \\
    & \hspace*{1cm} \leq \Bigg(\sum_{\xi\in I_N}\Big|\logE_{n\in[N]}\widetilde{a}(\Omega(n))
    \Big(\logE_{p\in\Pc_N}\f(n+p) - \logE_{m\in[N]}\f(m)\Big)\Big|^2\Bigg)^\frac{1}{2} \notag \\
    & \hspace*{1cm} \leq \Bigg(\sum_{\xi\in I_N}\logE_{n\in[N]}
    \Big|\logE_{p\in\Pc_N}\f(n+p) - \logE_{m\in[N]}\f(m)\Big|^2\Bigg)^\frac{1}{2}. 
\end{align}
The result follows from combining \eqref{applying_Tao}, \eqref{after_fourier_exp_eq} and \eqref{bounding_main_term_eq1}.
\end{proof}



%SECTION
\section{Small frequencies}\label{small_xi_section}

The goal of this section is to prove \cref{small_xi_thm}, which we first reduce to the following result.

\begin{theorem}\label{small_xi_Cesaro_thm}
Let $200\leq A_0 \leq \sqrt{\log\log{N}}$ and for $\xi\in I_N$, we define 
$$g(n) = g_{\xi,N,A_0}(n) = e\bigg(\frac{\xi\Omega(n)}{A_0\sqrt{\log\log{N}}}\bigg).$$
Let $B,H_1\geq 1$ satisfying $2\pi B\sqrt{\log\log{H_1}}\leq A_0\sqrt{\log\log{N}}$, and $\Pc\subset\P\cap[1,H_1]$.
Then for any $N^{\frac{1}{\log\log{N}}}\leq M\leq N^{2\log\log{N}}$, we have
\begin{multline*}
    \sum_{|\xi|\leq B} \Esub{n\in[N]}\Big|\logE_{p\in\Pc}g(n+p)-\Esub{m\in[M]}g(m)\Big|^2 \\
    \ll
    \frac{B^2}{A_0\sqrt{(\log\log{H_1})(\log\log{N})}} + \frac{BH_1}{N} + \frac{B(\log\log\log N)}{A_0\sqrt{\log\log{N}}}.
\end{multline*}
\end{theorem}

\begin{proof}[Proof of \cref{small_xi_Cesaro_thm} $\Longrightarrow$ \cref{small_xi_thm}]
    We let $\epsilon=(\log\log{N})^{-1}$.
    Using \cref{Cesaro_to_log}, Cauchy-Schwarz inequality and noting that for $M\in(N^\epsilon,N)$ it holds $\log\log{M}\asymp\log\log{N}$, we have 
    \begin{align}\label{M,M'_eq1}
        & \sum_{|\xi|\leq A^{5/2}}\logE_{n\in[N]}\Big|\logE_{p\in\Pc_N} \f(n+p) - \logE_{n\in[N]} \f(m)\Big|^2 \notag \\
        & = \sum_{|\xi|\leq A^{5/2}}\logE_{M\in(N^\epsilon,N)}\Esub{n\in[M]}\bigg|\logE_{p\in\Pc_N} \f(n+p) - \logE_{M'\in(N^\epsilon,N)}\Esub{m\in[M']}\f(m)\bigg|^2 \notag \\
        & \hspace*{.5cm} + \Oh\bigg(\frac{B}{\log\log{N}}\bigg) \notag \\
        & \leq \sup_{N^\epsilon<M,M'<N}\sum_{|\xi|\leq A^{5/2}} \Esub{n\in[M]}\Big|\logE_{p\in\Pc_N} \f(n+p) - \Esub{m\in[M']} \f(m) \Big|^2 + \Oh\bigg(\frac{B}{\log\log{N}}\bigg).
    \end{align}
    For $M,M'\in(N^\epsilon,N)$, we want to apply \cref{small_xi_Cesaro_thm} (with $M$ in the place of $N$ and $M'$ in the place of $M$) to bound the last sum. We choose $A_0 = A$, $B=A^{5/2}$, $H_1 = \exp\big((\log{M})^\frac{2}{(\log\log{M})^{4/9}}\big)$ and $\Pc=\Pc_N$. We check that the assumptions of \cref{small_xi_Cesaro_thm} are satisfied:
    \begin{itemize}
        \item $M^\frac{1}{\log\log{M}}\leq N^\frac{1}{\log\log{N}}\leq M' \leq N\leq M^{2\log\log{M}}$.
        \item We have 
        \begin{equation}\label{assumption_bound}
            \frac{B^2(\log\log{H_1})}{A_0^2(\log\log{M})}
            = \frac{(\log\log{N})^\frac{1}{3}(\log\log{M})^\frac{5}{9}}{\log\log{M}}
            \ll \frac{1}{(\log\log{N})^\frac{1}{9}},
        \end{equation}
        hence for $N$ sufficiently large, $2\pi B\sqrt{\log\log{H_1}}\leq A_0\sqrt{\log\log{M}}$.
        \item Since $\log{N}\leq 2(\log\log{M})(\log{M})\leq (\log{M})^2$ and $M\leq N$, we have 
        $$H = \exp\Big((\log{N})^\frac{1}{(\log\log{N})^{4/9}}\Big)
        \leq \exp\Big((\log{M})^\frac{2}{(\log\log{M})^{4/9}}\Big) = H_1,$$
        so $\Pc_N\subset[1,H_1]$.
    \end{itemize}
    Note that for this choice of $A_0$, $g(n) = \f(n)$, so for any such $M,M'$, by \cref{small_xi_Cesaro_thm} and since $\log\log{M}\asymp\log\log{N}$, we have
    \begin{equation}\label{M,M'_eq2}
        \sum_{|\xi|\leq B} \Esub{n\in[M]}  \Big|\logE_{p\in\Pc_N} \f(n+p) - \Esub{m\in[M']} \f(m) \Big|^2
        \ll \frac{\log\log\log{N}}{(\log\log{N})^\frac{1}{3}}.
    \end{equation}
    Combining \eqref{M,M'_eq1} with \eqref{M,M'_eq2} the result follows.
\end{proof}

\begin{proof}[Proof of \cref{small_xi_Cesaro_thm}]
We let $A_0,g,B,H_1,\Pc,M$ be as in the assumptions of \cref{small_xi_Cesaro_thm} and for simplicity we define 
$$S = \sum_{|\xi|\leq B} \Esub{n\in[N]}\Big|\logE_{p\in\Pc}g(n+p)-\Esub{m\in[M]}g(m)\Big|^2.$$

By expanding the square, we have
\begin{align*}
S & = \sum_{|\xi|\leq B}\Esub{n\in[N]}\bigg(\logE_{p,q\in\Pc}g(n+p)\overline{g(n+q)} 
\\
& \hspace*{3.2cm} - 2\Re\Big(\Big(\logE_{p\in\Pc}\overline{g(n+p)}\Big)\Big(\Esub{m\in[M]}g(m)\Big)\Big) + \Esub{m,\ell\in[M]}g(m)\overline{g(\ell)}\bigg) 
\\
& = \sum_{|\xi|\leq B}\bigg(
\logE_{p,q\in\Pc} \Esub{n\in[N]}g(n+p)\overline{g(n+q)} 
\\
& \hspace*{3.2cm} - 2\Re\bigg(\Esub{\substack{m\in[M] \\ n\in[N]}}\overline{g(n)}g(m)\bigg) + \Esub{m,\ell\in[M]}g(m)\overline{g(\ell)}\bigg) 
+ \Oh\bigg(\frac{BH_1}{N}\bigg) 
\\
& = \sum_{|\xi|\leq B}\Bigg(
\logE_{\substack{p,q\in\Pc \\ p\neq q}} \Esub{n\in[N]}g(n+p)\overline{g(n+q)}
\\
& \hspace*{3.2cm} - 2\Re\bigg(\Esub{\substack{m\in[M] \\ n\in[N]}}\overline{g(n)}g(m)\bigg) + \Esub{m,\ell\in[M]}g(m)\overline{g(\ell)}\Bigg) 
+ \Oh\bigg(\frac{B}{\log H_1}+\frac{BH_1}{N}\bigg) 
\\
& = \sum_{|\xi|\leq B}\Re\Bigg(
\logE_{\substack{p,q\in\Pc \\ p\neq q}} \Esub{n\in[N]}e\bigg(\frac{\xi \left(\Omega(n+p)-\Omega(n+q)\right)}{A_0\sqrt{\log\log{N}}}\bigg) 
\\
& \hspace*{0.2cm} - 2\Esub{\substack{m\in[M] \\ n\in[N]}}e\bigg(\frac{\xi\left(\Omega(m)-\Omega(n)\right)}{A_0\sqrt{\log\log{N}}}\bigg) + \Esub{m,\ell\in[M]} e\bigg(\frac{\xi \left(\Omega(m)-\Omega(\ell)\right)}{A_0 \sqrt{\log\log{N}}}\bigg)\Bigg) + \Oh\bigg(\frac{B}{\log H_1}+\frac{BH_1}{N}\bigg) \\
& \leq \sum_{|\xi|\leq B}\Bigg|
\logE_{\substack{p,q\in\Pc \\ p\neq q}} \Esub{n\in[N]}e\bigg(\frac{\xi \left(\Omega(n+p)-\Omega(n+q)\right)}{A_0\sqrt{\log\log{N}}}\bigg) 
\\
& \hspace*{0.2cm} - 2\Esub{\substack{m\in[M] \\ n\in[N]}}e\bigg(\frac{\xi\left(\Omega(m)-\Omega(n)\right)}{A_0\sqrt{\log\log{N}}}\bigg) + \Esub{m,\ell\in[M]} e\bigg(\frac{\xi \left(\Omega(m)-\Omega(\ell)\right)}{A_0 \sqrt{\log\log{N}}}\bigg)\Bigg|
+ \Oh\bigg(\frac{B}{\log H_1}+\frac{BH_1}{N}\bigg).
\end{align*}
Let $t_N$ and $\omega_N$ defined in \eqref{omega_N_def} respectively. Then, in view of \cref{lem_approx_Omega}, for all $|\xi|\leq B$, we have
\begin{multline*}
    \Esub{n\in[N]}e\bigg(\frac{\xi \left(\Omega(n+p)-\Omega(n+q)\right)}{A_0 \sqrt{\log\log{N}}}\bigg) \\
    = \Esub{n\in[N]}e\bigg(\frac{\xi \left(\omega_N(n+p)-\omega_N(n+q)\right)}{A_0 \sqrt{\log\log{N}}}\bigg)+\Oh\bigg(\frac{B(\log\log\log N)}{A_0\sqrt{\log\log{N}}}+\frac{BH_1}{A_0N}\bigg),
\end{multline*}
$$
\Esub{\substack{m\in[M] \\ n\in[N]}}e\bigg(\frac{\xi \left(\Omega(m)-\Omega(n)\right)}{A_0\sqrt{\log\log{N}}}\bigg)
=
\Esub{\substack{m\in[M] \\ n\in[N]}}e\bigg(\frac{\xi \left(\omega_N(m)-\omega_N(n)\right)}{A_0\sqrt{\log\log{N}}}\bigg)+\Oh\bigg(\frac{B(\log\log\log N)}{A_0\sqrt{\log\log{N}}}\bigg)$$
and
\[
\Esub{m,\ell\in[M]}e\bigg(\frac{\xi \left(\Omega(m)-\Omega(\ell)\right)}{A_0\sqrt{\log\log{N}}}\bigg)
=
\Esub{m,\ell\in[M]}e\bigg(\frac{\xi \left(\omega_N(m)-\omega_N(\ell)\right)}{A_0\sqrt{\log\log{N}}}\bigg)+\Oh\bigg(\frac{B(\log\log\log N)}{A_0\sqrt{\log\log{N}}}\bigg).
\]
We define 
\begin{align*}
    & S' = \sum_{|\xi|\leq B}\Bigg|\logE_{\substack{p,q\in\Pc \\ p\neq q}} \Esub{n\in[N]}e\bigg(\frac{\xi \left(\omega_N(n+p)-\omega_N(n+q)\right)}{A_0\sqrt{\log\log{N}}}\bigg) \\
    & \hspace*{2cm}
    - 2\Esub{\substack{m\in[M] \\ n\in[N]}}e\bigg(\frac{\xi\left(\omega_N(m)-\omega_N(n)\right)}{A_0\sqrt{\log\log{N}}}\bigg) + \Esub{m,\ell\in[M]} e\bigg(\frac{\xi \left(\omega_N(m)-\omega_N(\ell)\right)}{A_0 \sqrt{\log\log{N}}}\bigg)\Bigg|
\end{align*}
and then we have
\begin{equation}
\label{eqn_3to4}
S 
\leq S' + \Oh\bigg(\frac{B}{\log{H_1}} + \frac{BH_1}{N} + \frac{B(\log\log\log N)}{A_0\sqrt{\log\log{N}}}\bigg).
\end{equation}

Let
$K = \lfloor(\log\log{N})^5\rfloor$ and
$e_K(x)=\sum_{k=0}^K \frac{(2\pi i)^k x^k}{k!}$
be the degree $K$ approximation of $e(x)$.
Since $B\leq \sqrt{\log\log{N}}$ and $A_0\geq 200$, we have $K\geq \frac{200 B (\log\log{N})^4}{A_0}$. So using the definition of $\omega_N$, \cref{lem_approx_deg_K} and that $\log\log{M}\asymp\log\log{N}$, we have

\begin{align}\label{eqn_S'}
S' &= \sum_{|\xi|\leq B}\Bigg|
\logE_{\substack{p,q\in\Pc \\ p\neq q}} \Esub{n\in[N]}e\bigg(\frac{\xi \big(\sum_{r\leq t_N}\1_{r\mid n+p}-1_{r\mid n+q}\big)}{A_0\sqrt{\log\log{N}}}\bigg) \notag \\
& \hspace*{1.5cm} - 2\Esub{\substack{m\in[M] \\ n\in[N]}}e\bigg(\frac{\xi\big(\sum_{r\leq t_N}\1_{r\mid m}-\1_{r\mid n}\big)}{A_0\sqrt{\log\log{N}}}\bigg) + \Esub{m,\ell\in[M]} e\bigg(\frac{\xi\big(\sum_{r\leq t_N}\1_{r\mid m}-\1_{r\mid \ell}\big)}{A_0\sqrt{\log\log{N}}}\bigg)\Bigg|
\notag \\
&=\sum_{|\xi|\leq B}\Bigg|
\logE_{\substack{p,q\in\Pc \\ p\neq q}} \Esub{n\in[N]}e_K\bigg(\frac{\xi \big(\sum_{r\leq t_N}\1_{r\mid n+p}-1_{r\mid n+q}\big)}{A_0\sqrt{\log\log{N}}}\bigg) \notag \\
& \hspace*{1.5cm} - 2\Esub{\substack{m\in[M] \\ n\in[N]}}e_K\bigg(\frac{\xi\big(\sum_{r\leq t_N}\1_{r\mid m}-\1_{r\mid n}\big)}{A_0\sqrt{\log\log{N}}}\bigg) + \Esub{m,\ell\in[M]} e_K\bigg(\frac{\xi\big(\sum_{r\leq t_N}\1_{r\mid m}-\1_{r\mid \ell}\big)}{A_0\sqrt{\log\log{N}}}\bigg)\Bigg|
\notag \\
&\hspace*{8.9cm} + \Oh\bigg(B e^{-K}+\frac{B}{(\log\log N)^4}+\frac{BH_1}{N}\bigg) \notag \\
&= \sum_{|\xi|\leq B}\Bigg|
\logE_{\substack{p,q\in\Pc \\ p\neq q}} \Esub{n\in[N]}
\sum_{k=0}^K \frac{(2\pi i)^k \xi^k}{A_0^k k!(\log\log N)^{\frac{k}{2}}}
\bigg(\sum_{r\leq t_N}\1_{r\mid n+p}-\1_{r\mid n+q}\bigg)^k
\notag \\
& \hspace*{1.5cm}
-2\Esub{\substack{m\in[M] \\ n\in[N]}}
\sum_{k=0}^K \frac{(2\pi i)^k \xi^k}{A_0^k k!(\log\log N)^{\frac{k}{2}}}\bigg(\sum_{r\leq t_N} \1_{r\mid m}-\1_{r\mid n}\bigg)^k
\notag \\
& \hspace*{1.5cm}
+ \Esub{m,\ell\in[M]} \sum_{k=0}^K \frac{(2\pi i)^k \xi^k}{A_0^k k!(\log\log N)^{\frac{k}{2}}} \bigg(\sum_{r\leq t_N}\1_{r\mid m}-\1_{r\mid \ell}\bigg)^k\Bigg|
+\Oh\bigg(\frac{B}{(\log\log N)^4}+\frac{BH_1}{N}\bigg)
\notag \\
& \leq 
\sum_{|\xi|\leq B}\sum_{k=1}^K \frac{(2\pi\xi)^k}{A_0^k k! (\log\log N)^{\frac{k}{2}}}
\Bigg|
\logE_{\substack{p,q\in\Pc \\ p\neq q}} \Esub{n\in[N]} \bigg(\sum_{r\leq t_N}\1_{r\mid n+p}-\1_{r\mid n+q}\bigg)^k
\notag \\
& \hspace*{4.3cm}
- 2\Esub{\substack{m\in[M] \\ n\in[N]}}
\bigg(\sum_{r\leq t_N} \1_{r\mid m}-\1_{r\mid n}\bigg)^k
+\Esub{m,\ell\in[M]}\bigg(\sum_{r\leq t_N}\1_{r\mid m}-\1_{r\mid \ell}\bigg)^k\Bigg| \notag \\
& \hspace*{10.5cm} + \Oh\bigg(\frac{B}{(\log\log N)^4}+\frac{BH_1}{N}\bigg)
\notag \\
& = \sum_{|\xi|\leq B}\sum_{k=1}^K \frac{(2\pi\xi)^k}{A_0^k k! (\log\log N)^{\frac{k}{2}}}
S''(k) + \Oh\bigg(\frac{B}{(\log\log N)^4}+\frac{BH_1}{N}\bigg) \notag \\
& = \sum_{k=1}^K \frac{(2\pi)^k}{A_0^k k! (\log\log{N})^\frac{k}{2}}S''(k) \bigg(\underbrace{\sum_{|\xi|\leq B}\xi^k}_{=\frac{B^{k+1}-1}{k-1}}\bigg) + \Oh\bigg(\frac{B}{(\log\log N)^4}+\frac{BH_1}{N}\bigg)
\notag \\
& \ll \sum_{k=1}^K \frac{(2\pi)^k B^{k+1}}{A_0^{k} k k!(\log\log{N})^\frac{k}{2}}S''(k) + \frac{B}{(\log\log{N})^4}+\frac{BH_1}{N},
\end{align}
where we have set
\begin{multline}\label{eqn_S''_defn}
    S''(k)
    =
    \Bigg|\logE_{\substack{p,q\in\Pc \\ p\neq q}} \Esub{n\in[N]} \bigg(\sum_{r\leq t_N}\1_{r\mid n+p}-\1_{r\mid n+q}\bigg)^k \\
    - 2\Esub{\substack{m\in[M] \\ n\in[N]}}\bigg(\sum_{r\leq t_N} \1_{r\mid m}-\1_{r\mid n}\bigg)^k
    +\Esub{m,\ell\in[M]}\bigg(\sum_{r\leq t_N}\1_{r\mid m}-\1_{r\mid \ell}\bigg)^k\Bigg|,
\end{multline}
and the sum over $r$ is a sum over primes. Throughout, the variable $r$ will denote a prime.

Given $u_1,\ldots,u_j\geq 0$ with $u_1+\ldots+u_j=k$, let
\[
\binom{k}{u_1,\ldots,u_j}
\]
denote the multinomial coefficient, which expresses the number of unique permutations of a word of length $k$ using $j$ distinct letters such that the $i$-th letter occurs $u_i$-any times. 

By expanding, we have
\begin{align*}
\bigg(\sum_{r\leq t_N}\1_{r\mid n+p} & - \1_{r\mid n+q}\bigg)^k
= \sum_{r_1,\ldots,r_k\leq t_N} \prod_{i=1}^k \left(\1_{r_i\mid n+p} - \1_{r_i\mid n+q}\right) \\
& = \sum_{j=1}^k \sum_{\substack{u_1,\ldots,u_j\geq 1 \\ u_1+\ldots+u_j=k}}\binom{k}{u_1,\ldots,u_j}\sum_{\substack{r_1,\ldots,r_j\leq t_N \\ \text{pairwise distinct}}}\prod_{i=1}^j \left(\1_{r_i\mid n+p} - \1_{r_i\mid n+q}\right)^{u_i} \\
& = 
\sum_{j=1}^k \sum_{\substack{u_1,\ldots,u_j\geq 1 \\ u_1+\ldots+u_j=k}}\binom{k}{u_1,\ldots,u_j}\sum_{\substack{r_1,\ldots,r_j\leq t_N \\ \text{pairwise distinct}}}\prod_{i=1}^j \left(\1_{r_i\mid n+p} - \1_{r_i\mid n+q}\right)^{[u_i]},
\end{align*}
where $[u]_2=1$ if $u$ is odd and $[u]_2=2$ if $u$ is even.
In view of \cref{lem_independence_periodic_fctns_np_1}, for any pairwise distinct primes $r_1,\ldots,r_j$, letting $\epsilon = (\log\log{N})^{-1}$, we have
$$\Esub{n\in[N]}\prod_{i=1}^{j}\left(\1_{r_i\mid n+p}-\1_{r_i\mid n+q}\right)^{[u_i]_2} 
= \prod_{i=1}^{j}\Esub{n\in[N]}\left(\1_{r_i\mid n+p}-\1_{r_i\mid n+q}\right)^{[u_i]_2}
+ \Oh\bigg(\frac{2^{j^2}}{N}\bigg),$$
$$\Esub{\substack{m\in[M] \\ n\in[N]}} \prod_{i=1}^{j} \left(\1_{r_i\mid m}-\1_{r_i\mid n}\right)^{[u_i]_2}
= \prod_{i=1}^{j}\Esub{\substack{m\in[M] \\ n\in[N]}} \left(\1_{r_i\mid m}-\1_{r_i\mid n}\right)^{[u_i]_2}
+ \Oh\bigg(\frac{2^{j^2}}{N^\epsilon}\bigg)$$
and 
$$\Esub{m,\ell\in[M]} \prod_{i=1}^{j} \left(\1_{r_i\mid m}-\1_{r_i\mid \ell}\right)^{[u_i]_2}
= \prod_{i=1}^{j}\Esub{m,\ell\in[M]}\left(\1_{r_i\mid m}-\1_{r_i\mid \ell}\right)^{[u_i]_2}
+ \Oh\bigg(\frac{2^{j^2}}{N^\epsilon}\bigg).$$
It follows that
\begin{align*}
&\Esub{n\in[N]}\bigg(\sum_{r\leq t_N}\1_{r\mid n+p}-\1_{r\mid n+q}\bigg)^k
\\
& \hspace*{1.1cm} = \sum_{j=1}^k \sum_{\substack{u_1,\ldots,u_j\geq 1 \\ u_1+\ldots+u_j=k}} \binom{k}{u_1,\ldots,u_{t_N}} \sum_{\substack{r_1,\ldots,r_j\leq t_N \\ \text{pairwise distinct}}}\Bigg(\prod_{i=1}^j \Esub{n\in[N]}  
\left(\1_{r_i\mid n+p}-\1_{r_i\mid n+q}\right)^{[u_i]_2} \\
& \hspace*{12.7cm} + \Oh\bigg(\frac{2^{j^2}}{N^\epsilon}\bigg)\Bigg)
\\
&\hspace*{1.1cm}
= \sum_{j=1}^k \sum_{\substack{u_1,\ldots,u_{j}\geq 1 \\ u_1+\ldots+u_j=k}} \binom{k}{u_1,\ldots,u_j}\underbrace{\sum_{\substack{r_1,\ldots,r_j\leq t_N \\ \text{pairwise distinct}}}\prod_{i=1}^{j} \Esub{n\in[N]}  
\left(\1_{r_i\mid n+p}-\1_{r_i\mid n+q}\right)^{[u_i]_2}}_{:=\Phi_1(u_1,\ldots,u_j;p,q)} \\
& \hspace*{12.7cm} + \Oh\bigg(\frac{2^{k^2}t_N^k}{N^\epsilon}\bigg),
\end{align*}
where for the calculation of the error term in the last step above we used that
\begin{align*}
    \sum_{j=1}^k \sum_{\substack{u_1,\ldots,u_j\geq 1 \\ u_1+\ldots+u_j=k}} \binom{k}{u_1,\ldots,u_{t_N}} \sum_{\substack{r_1,\ldots,r_j\leq t_N \\ \text{pairwise distinct}}} 2^{j^2}
    & \leq 2^{k^2} \sum_{j=1}^k \binom{t_N}{j} \sum_{\substack{u_1,\ldots,u_j\geq 1 \\ u_1+\ldots+u_j=k}} \binom{k}{u_1,\ldots,u_{t_N}} \\
    & = 2^{k^2}\sum_{\substack{u_1,\ldots,u_{t_N}\geq 0 \\ u_1+\dots +u_{t_N}=k}}\binom{k}{u_1,\ldots,u_{t_N}}
    = 2^{k^2}t_N^k,
\end{align*}
where the last equality follows from the multinomial theorem.
Similar calculations show 
\begin{multline*}
\Esub{\substack{m\in[M] \\ n\in[N]}}\bigg(\sum_{r\leq t_N}\1_{r\mid m}-\1_{r\mid n}\bigg)^k
\\
=
\sum_{j=1}^k\sum_{\substack{u_1,\ldots,u_j\geq 1 \\ u_1+\ldots+u_j=k}} \binom{k}{u_1,\ldots,u_j} \underbrace{\sum_{\substack{r_1,\ldots,r_j\leq t_N \\ \text{pairwise distinct}}}\prod_{i=1}^j\Esub{\substack{m\in[M] \\ n\in[N]}} 
\left(\1_{r_i\mid m}-\1_{r_i\mid n}\right)^{[u_i]_2}}_{:=\Phi_2(u_1,\ldots,u_j)} + \Oh\bigg(\frac{2^{k^2}t_N^k}{N^\epsilon}\bigg)
\end{multline*}
and
\begin{multline*}
\Esub{m,\ell\in[M]}\bigg(\sum_{r\leq t_N}\1_{r\mid m}-\1_{r\mid \ell}\bigg)^k
\\
=
\sum_{j=1}^k\sum_{\substack{u_1,\ldots,u_j\geq 1 \\ u_1+\ldots+u_j=k}} \binom{k}{u_1,\ldots,u_j} \underbrace{\sum_{\substack{r_1,\ldots,r_j\leq t_N \\ \text{pairwise distinct}}}\prod_{i=1}^j \Esub{m,\ell\in[M]}  
\left(\1_{r_i\mid m}-\1_{r_i\mid \ell}\right)^{[u_i]_2}}_{:=\Phi_3(u_1,\ldots,u_j)} + \Oh\bigg(\frac{2^{k^2}t_N^k}{N^\epsilon}\bigg).
\end{multline*}
Plugging the last three estimates in \eqref{eqn_S''_defn} and using the triangle inequality gives 
\begin{align}\label{eqn_S''_2}
    S''(k) \leq \sum_{j=1}^k\sum_{\substack{u_1,\ldots,u_j\geq 1 \\ u_1+\ldots+u_j=k}} \binom{k}{u_1,\ldots,u_j}\bigg| \logE_{\substack{p,q\in\Pc \\ p\neq q}} \Phi_1(&u_1,\ldots,u_j;p,q) - 2\Phi_2(u_1,\ldots,u_j) \notag \\
    & + \Phi_3(u_1,\ldots,u_j)\bigg| + \Oh\bigg(\frac{2^{k^2}t_N^k}{N^\epsilon}\bigg).
\end{align}

For $u_1,\ldots,u_j\geq 1$ with $u_1+\ldots+u_j = k$,
we distinguish two cases:
\begin{itemize}
    \item If $u_i$ is odd for some $1\leq i\leq j$, then for any $r_1,\ldots,r_j\leq t_N$ pairwise distinct primes, we have
    $$\bigg|\Esub{n\in[N]}\1_{r_i\mid n+p}-\1_{r_i\mid n+q}\bigg| < \frac{2}{N}, \quad
    \bigg|\Esub{\substack{m\in[M] \\ n\in[N]}}\1_{r_i\mid m}-\1_{r_i\mid n}\bigg|<\frac{1}{N^\epsilon}$$
    and 
    $$\bigg|\Esub{m,\ell\in[M]}\1_{r_i\mid m}-\1_{r_i\mid \ell}\bigg|<\frac{1}{N^\epsilon}.$$
    Then
    \begin{multline}\label{Phi_bound1}
        \bigg|\logE_{\substack{p,q\in\Pc \\ p\neq q}}\Phi_1(u_1,\ldots,u_j;p,q)-2\Phi_2(u_1,\ldots,u_j)+\Phi_3(u_1,\ldots,u_j)\bigg| \\
        \ll \frac{1}{N^\epsilon}\sum_{\substack{r_1,\ldots,r_j\leq t_N \\ \text{pairwise distinct}}}1
        \leq \frac{1}{N^\epsilon}\binom{t_N}{j}. 
    \end{multline}
    The contribution of these terms in \eqref{eqn_S''_2} is 
    \begin{equation}\label{eqn_1contr_S''}
        \ll \frac{1}{M}\sum_{j=1}^k \binom{t_N}{j} \sum_{\substack{u_1,\ldots,u_j\geq 1 \\ u_1+\ldots+u_j=k}} \binom{k}{u_1,\ldots,u_{t_N}} 
        = \frac{1}{N^\epsilon} \sum_{\substack{u_1,\ldots,u_{t_N}\geq 0 \\ u_1+\ldots +u_{t_N}=k}}\binom{k}{u_1,\ldots,u_{t_N}} 
        = \frac{t_N^k}{N^\epsilon},
    \end{equation}
where the last equality follows again from the multinomial theorem.
\item On the other hand, if $u_i$ is even for all $1\leq i\leq j$, then any $r_1,\ldots,r_j\leq t_N$ pairwise distinct primes, we have
    \begin{align*}
        \Esub{n\in[N]}\big(\1_{r_i\mid n+p}-\1_{r_i\mid n+q}\big)^{[u_i]_2}
        & = \frac{2}{r_i}\big(1-\1_{r_i\mid p-q}\big)+\Oh\bigg(\frac{1}{N}\bigg), \\
        \Esub{\substack{m\in[M] \\ n\in[N]}}\big(\1_{r_i\mid m} - \1_{r_i\mid n}\big)^{[u_i]_2} 
        & = \frac{2}{r_i} + \Oh\bigg(\frac{1}{N^\epsilon}\bigg), \\
        \Esub{m,\ell\in[M]}\big(\1_{r_i\mid m}-\1_{r_i\mid \ell}\big)^{[u_i]_2}
        & = \frac{2}{r_i}+\Oh\bigg(\frac{1}{N^\epsilon}\bigg).
    \end{align*}
    Taking product over $1\leq i \leq j$ yields
    \begin{align*}
        \prod_{i=1}^j \Esub{n\in[N]}\big(\1_{r_i\mid n+p}-\1_{r_i\mid n+q}\big)^{[u_i]_2}
        & = \prod_{i=1}^j \frac{2}{r_i}\big(1-\1_{r_i\mid p-q}\big) + \Oh\Bigg(\sum_{i=1}^j \frac{1}{N^i}\binom{j}{j-i}\Bigg) \\
        & = \prod_{i=1}^j \frac{2}{r_i}\big(1-\1_{r_i\mid p-q}\big) + \Oh\bigg(\frac{j^2}{N}\bigg)
    \end{align*}
    and similar estimates hold for the other two terms. Thus, we obtain
    \begin{align*}
        & \bigg|\logE_{\substack{p,q\in\Pc \\ p\neq q}} 
        \Phi_1(u_1,\ldots,u_j;p,q)-2\Phi_2(u_1,\ldots,u_j) + \Phi_3(u_1,\ldots,u_j)\bigg| =
        \\
        & \hspace*{1.5cm} = \Bigg|\logE_{\substack{p,q\in\Pc \\ p\neq q}}
        \sum_{\substack{r_1,\ldots,r_j\leq t_N \\ \text{pairwise distinct}}}\Bigg(
        \prod_{i=1}^{j}\frac{2}{r_i}\left(1-\1_{r_i\mid p-q}\right) - 
        \prod_{i=1}^{j}\frac{2}{r_i}\Bigg)\Bigg|
        + \Oh\Bigg(\frac{j^2\binom{t_N}{j}}{N^\epsilon}\Bigg)
        \\
        & \hspace*{1.5cm} = \logE_{\substack{p,q\in\Pc \\ p\neq q}}\sum_{\substack{r_1,\ldots,r_j\leq t_N \\ \text{pairwise distinct}}}\prod_{i=1}^j\frac{2}{r_i}\1_{r_i\mid p-q} + \Oh\Bigg(\frac{j^2\binom{t_N}{j}}{N^\epsilon}\Bigg) \\
        & \hspace*{1.5cm} \leq j \Bigg(\sum_{r_1\leq H_1}\frac{2}{r_1}\bigg(\underbrace{\logE_{\substack{p,q\in\Pc \\ p\neq q}}\1_{r_1\mid p-q}}_{\approx\frac{1}{r_1}}\bigg)\Bigg) \Bigg(\sum_{r_2,\ldots,r_j\leq H}\frac{2}{r_2\cdots r_j}\Bigg) + \Oh\Bigg(\frac{j^2\binom{t_N}{j}}{N^\epsilon}\Bigg) \\
        & \hspace*{1.5cm} \ll j \Bigg(\underbrace{\sum_{r_1\leq H_1}\frac{2}{r_1^2}}_{\ll 1}\Bigg) \Bigg(\underbrace{\sum_{r_2,\ldots r_j\leq H_1}\frac{2}{r_2\cdots r_j}}_{\ll (\log\log{H_1})^{j-1}}\Bigg) + \frac{j^2\binom{t_N}{j}}{N^\epsilon}
        \ll j(\log\log{H_1})^{j-1} + \frac{j^2\binom{t_N}{j}}{N^\epsilon}.
    \end{align*}
    Since $u_1,\ldots,u_j$ are even, notice that $j\leq \frac{k}{2}$. Then the contribution of these terms in \eqref{eqn_S''_2} is 
    \begin{align}\label{eqn_2contr_S''}
        & \ll \sum_{j=1}^{\frac{k}{2}} \Bigg(j(\log\log{H_1})^{j-1}+ \frac{j^2\binom{t_N}{j}}{N^\epsilon}\Bigg) \sum_{\substack{u_1,\ldots,u_j\geq 1 \notag \\ u_1+\ldots+u_j = \left\lfloor\frac{k}{2}\right\rfloor}}\binom{\left\lfloor\frac{k}{2}\right\rfloor}{u_1,\ldots,u_j} \\
        & \leq k(\log\log{H_1})^{\frac{k}{2}-1}\sum_{j=1}^{\frac{k}{2}}\sum_{\substack{u_1,\ldots,u_j\geq 1 \\ u_1+\ldots+u_j = \left\lfloor\frac{k}{2}\right\rfloor}}\binom{\left\lfloor\frac{k}{2}\right\rfloor}{u_1,\ldots,u_j} + \frac{k^2}{N^\epsilon} \sum_{\substack{u_1,\ldots,u_{t_N}\geq 0 \\ u_1+\ldots+u_{t_N} = k}}\binom{k}{u_1,\ldots,u_{t_N}}\notag \\
        & \ll k \bigg(\frac{k}{2}\bigg)^{\frac{k}{2}} (\log\log{H_1})^{\frac{k}{2}-1} + \frac{k^2 t_N^k}{N^\epsilon}.
    \end{align}
\end{itemize}

Substituting the bounds \eqref{eqn_1contr_S''} and \eqref{eqn_2contr_S''} in \eqref{eqn_S''_2}, we obtain
$$S''(k) \ll \frac{2^{k^2}t_N^k}{N^\epsilon} + k \bigg(\frac{k}{2}\bigg)^{\frac{k}{2}}(\log\log{H_1})^{\frac{k}{2}-1}.$$
Substituting the above in \eqref{eqn_S'} we get
\begin{align}\label{S'_bound1}
    S' & \ll \sum_{k=1}^K \frac{(2\pi)^k B^{k+1}}{A_0^k k k! (\log\log{N})^\frac{k}{2}}\Bigg(\frac{2^{k^2}t_N^k}{N^\epsilon} + k \bigg(\frac{k}{2}\bigg)^{\frac{k}{2}} (\log\log{H_1})^{\frac{k}{2}-1}\Bigg) + \frac{B}{(\log\log{N})^4} + \frac{BH_1}{N} \notag \\
    & \leq \frac{2^{K^2} t_N^K B}{N^\epsilon} \sum_{k=1}^K \frac{1}{k k!}\bigg(\frac{2\pi B}{A_0\sqrt{\log\log{N}}}\bigg)^k
    + \frac{2\pi B}{\log\log{H_1}} \sum_{k=1}^K \frac{k^\frac{k}{2}}{2^\frac{k}{2}k!}\bigg(\frac{B\sqrt{\log\log{H_1}}}{A_0\sqrt{\log\log{N}}}\bigg)^k \notag \\
    & \hspace*{9cm} + \frac{B}{(\log\log{N})^4} + \frac{BH_1}{N}.
\end{align}
By assumption, we have
$$\frac{2\pi B\sqrt{\log\log{H_1}}}{A_0\sqrt{\log\log{N}}}<1.$$
Moreover, by Stirling formula, we can see that 
$$\frac{k^\frac{k}{2}}{2^\frac{k}{2}k!}\ll \frac{\Big(\frac{e}{\sqrt{2}}\Big)^k}{k^\frac{k}{2}},$$ 
so we obtain
$$\sum_{k=1}^K \frac{1}{k k!}\bigg(\frac{2\pi B}{A_0\sqrt{\log\log{N}}}\bigg)^k
< \sum_{k=1}^K \frac{1}{k k!} \ll 1$$
and
$$
\sum_{k=1}^K \frac{k^\frac{k}{2}}{2^\frac{k}{2}k!}\bigg(\frac{B\sqrt{\log\log{H_1}}}{A_0\sqrt{\log\log{N}}}\bigg)^k
\leq \frac{B\sqrt{\log\log{H_1}}}{A_0\sqrt{\log\log{N}}} \sum_{k=1}^K \frac{\Big(\frac{e}{\sqrt{2}}\Big)^k}{k^\frac{k}{2}}
\ll \frac{B\sqrt{\log\log{H_1}}}{A_0\sqrt{\log\log{N}}}.$$
Plugging this bounds in \eqref{S'_bound1} we get
\begin{align}\label{S'_bound2}
    S' & \ll \frac{2^{K^2} t_N^K B}{N^\epsilon} + \frac{B^2}{A_0\sqrt{(\log\log{H_1})(\log\log{N})}} + \frac{BH_1}{N} \notag \\
    & \leq \frac{2^{(\log\log{N})^{10}}N^\frac{1}{(\log\log{N})^3}B}{N^\frac{1}{\log\log{N}}} + \frac{B^2}{A_0\sqrt{(\log\log{H_1})(\log\log{N})}} + \frac{BH_1}{N} \notag \\
    & \ll \frac{B^2}{A_0\sqrt{(\log\log{H_1})(\log\log{N})}} + \frac{BH_1}{N}.
\end{align}
Plugging \eqref{S'_bound2} in \eqref{eqn_3to4}, we have
$$S \ll \frac{B^2}{A_0\sqrt{(\log\log{H_1})(\log\log{N})}} + \frac{BH_1}{N} + \frac{B(\log\log\log N)}{A_0\sqrt{\log\log{N}}}.$$
\end{proof}



%SECTION
\section{Large frequencies}\label{large_xi_section}

The goal of this subsection is to prove \cref{large_xi_thm}.
Expanding the square in 
$$\sum_{\substack{\xi\in I_N \\ |\xi|>A^{5/2}}}
\logE_{n\in[N]}\Big|\logE_{p\in\Pc_N}\f(n+p) - \logE_{m\in[N]} \f(m)\Big|^2$$
we see that \cref{large_xi_thm} follows from the next two propositions.

\begin{proposition}\label{large_xi_thm1}
We have
    $$\sum_{\substack{\xi\in I_N \\ |\xi|>A^{5/2}}}\Big|\Esub{n\in[N]} \f(n)\Big| 
    \ll \exp\big(-(\log\log{N})^\frac{1}{3}\big).$$
\end{proposition}

\begin{proposition}\label{large_xi_thm2}
We have
    $$\sum_{\substack{\xi\in I_N \\ |\xi|>A^{5/2}}}
    \Esub{n\in[N]}\Big|\logE_{p\in\Pc_N}\f(n+p)\Big|^2 
    \ll \frac{1}{(\log\log{N})^\frac{1}{3}}.$$
\end{proposition}

We need these propositions for logarithmic averages over $n$, but we can prove them for \Cesaro{} averages, which, in view of \cref{Cesaro_to_log}, makes the results stronger.

To prove \cref{large_xi_thm1} we use \Halasz's mean value theorem, while the proof of \cref{large_xi_thm2} requires using a deeper result of \Matomaki-\Radziwill-Tao \cite{MRT15}. For both cases we need to introduce the notion of distance between two multiplicative functions.

\begin{defn}\label{defn_distance}
    Let $f,g\colon\N\to\C$ be two bounded multiplicative functions. The \textit{distance (up to $N$)} between $f$ and $g$ is defined as $$\Dist(f,g;N)=\Bigg(\sum_{p\leq N}\frac{1-\Re(f(p)\overline{g(p)})}{p}\Bigg)^{\frac{1}{2}}.$$ 
\end{defn}

Given a $1$-bounded multiplicative function $f$, we define 
$$M_0(f;N) := \inf_{|t|\leq N} \Dist(f, n\mapsto n^{it};N)^2.$$

\subsection{Proof of \cref{large_xi_thm1}}

We will use a version of \Halasz's mean value theorem due to Granville and Soundararajan.

\begin{theorem}[{\cite[Theorem 1.2]{granville2003decay}}]\label{Halasz_GS}
For any $1$-bounded multiplicative function $f$, we have 
$$\Big|\Esub{n\in[N]} f(n)\Big| \ll \exp\bigg(-\frac{1}{16}\inf_{|t|\leq\log{N}}\Dist(f,n\mapsto n^{it};N)^2\bigg).$$
\end{theorem}

In the next two lemmas we provide some lower bounds for the distance of $\f$ from Archimedean characters $n^{it}$, $t\in\R$.

\begin{lemma}\label{dist1}
For any $\xi\in I_N$ and $|t|\ll1$, we have 
\begin{equation}\label{dist1.eq1}
    \Dist(\f,n\mapsto n^{it};N)^2 
    = \bigg(1-\cos\bigg(\frac{2\pi\xi}{|I_N|}\bigg)\bigg)\log\log{N} 
    + \cos\bigg(\frac{2\pi\xi}{|I_N|}\bigg)\log(1+|t|\log{N}) + \Oh(1).
\end{equation}
In particular, for $N$ sufficiently large, we have 
\begin{equation}\label{dist1.eq2}
    \inf_{|t|\ll 1}\Dist(\f,n\mapsto n^{it};N)^2\geq \max\bigg\{\frac{\pi^2\xi^2}{2A^2}, \frac{\pi^2(\log\log{N})}{A}\bigg\} + \Oh(1).
\end{equation}
\end{lemma}

\begin{lemma}\label{dist2}
For any $\epsilon>0$, $B>0$, $\xi\in I_N$, and $2\leq |t|\leq N^B$, we have
\begin{equation}\label{dist2.eq1}
    \Dist(\f,n\mapsto n^{it}; N)^2\geq \bigg(\frac{1}{3}-\epsilon\bigg)\log\log{N}+\Oh_{B,\epsilon}(1).
\end{equation}
\end{lemma}

\begin{proof}[Proof of \cref{dist1}]
Let $s=1+\frac{1}{\log{N}}+it$, $|t|\ll1$, and for $\Re(s)>1$, define $$F(s)=\sum_{n\geq1}\frac{\f(n)}{n^s}.$$ Then we have
\begin{equation}\label{L-f}
    \log{F(s)}=\sum_{p}\frac{\f(p)}{p^s}+\Oh(1) = e\bigg(\frac{\xi}{|I_N|}\bigg)\log{\zeta(s)} + \Oh(1).
\end{equation}
Now, using the following classic estimate for the Riemann zeta function
$$\zeta(s)=\frac{s}{s-1}+O\bigg(\frac{|s|}{\sigma}\bigg)=\frac{1}{s-1}+\Oh(1),
\quad \sigma=\Re(s)>0,$$
we get that
\begin{align}\label{lz}
    \log{\zeta(s)} & =-\log\bigg(\frac{1}{\log{N}}+it\bigg)+\Oh(1) \\
    & = -\log\bigg|\frac{1}{\log{N}}+it\bigg| -  i\arg\bigg(\frac{1}{\log{N}}+it\bigg)+\Oh(1) \notag \\
    & = \log\log{N} - \log(1+|t|\log{N}) + \Oh(1). \notag
\end{align}
Combining \eqref{L-f} and \eqref{lz}, we obtain that
\begin{equation}\label{ReF}
    \Re(\log{F(s)})
    = \cos\bigg(\frac{2\pi\xi}{|I_N|}\bigg)(\log\log{N}-\log(1+|t|\log{N})) + \Oh(1)
\end{equation}
Now, using a classical asymptotic relating the distance of a multiplicative function from twisted characters with its Dirichlet L-function at $1+\frac{1}{\log{N}}+it$ (see \cite[Eq. (2.2)]{GS_large_char_sums}), we have
$$\exp(-\Dist(\f,n\mapsto n^{it};N)^2)=\frac{|F(s)|}{\log{N}} + \Oh(1).$$ 
Combining this with the fact that $\log|F(s)|=\Re(\log{F(s)})$, we see that
$$\Dist(\f,n\mapsto n^{it};N)^2=\log\log{N}-\log|F(s)|=\log\log{N}-\Re(\log{F(s)})$$ and then, by substituting \eqref{ReF} to the right-hand side of the above we obtain \eqref{dist1.eq1}.
\par
We will now prove the second part. We distinguish two cases depending on the size of $|\xi|$:
\begin{itemize}
    \item If $|\xi|\leq|I_N|/4$, then $\cos(2\pi\xi/|I_N|)\geq0$, hence \eqref{dist1.eq1} gives
    $$\Dist(\f,n\mapsto n^{it};N)^2 \geq \bigg(1-\cos\bigg(\frac{2\pi\xi}{|I_N|}\bigg)\bigg)\log\log{N} + O(1).$$
    $1-\cos(2\pi\xi/|I_N|)$ is increasing as a function of $|\xi|$ and it is also positive. Moreover, by the Laurent expansion of $1-\cos(1/x)$ at $x=\infty$, we have
    \begin{equation}\label{Laurent_exp}
        1-\cos\bigg(\frac{1}{x}\bigg) = \frac{1}{2x^2} + \Oh\bigg(\frac{1}{x^4}\bigg).
    \end{equation}
    \begin{itemize}
        \item  If $|\xi|\leq \sqrt{|I_N|}$, then \eqref{Laurent_exp} gives that
        for $N$ sufficiently large, 
        $$1-\cos\bigg(\frac{2\pi\xi}{|I_N|}\bigg) 
        = \frac{2\pi^2\xi^2}{|I_N|^2} + \Oh\bigg(\frac{1}{|I_N|^2}\bigg),$$
        hence 
        $$\Dist(\f,n\mapsto n^{it};N)^2 
        \geq \frac{\pi^2\xi^2}{2A^2} + \Oh(1).$$
        \item If $\sqrt{|I_N|}<|\xi|\leq|I_N|/4$, then by monotonicity and \eqref{Laurent_exp}, we have
        $$1-\cos\bigg(\frac{2\pi\xi}{|I_N|}\bigg) > 1 - \cos\bigg(\frac{2\pi}{\sqrt{|I_N|}}\bigg) = \frac{2\pi^2}{|I_N|} + \Oh\bigg(\frac{1}{|I_N|^2}\bigg),$$
        thus,
        $$\Dist(\f,n\mapsto n^{it};N)^2 \geq \frac{\pi^2(\log\log{N})}{A} + \Oh(1).$$
    \end{itemize}
    \item If $|I_N|/4\leq |\xi| \leq |I_N|/2$, then $\cos(2\pi\xi/|I_N|)\leq 0$. Rewriting \eqref{dist1.eq1}, we have
    $$\Dist(\f,n\mapsto n^{it};N)^2
    = \log\log{N} + \cos\bigg(\frac{2\pi\xi}{|I_N|}\bigg)\log\bigg(\frac{1+|t|\log{N}}{\log{N}}\bigg) + \Oh(1)$$
    and now we distinguish cases on the size of $|t|$:
    If $|t|\leq 1/2$, then the second term above is non-negative for $N$ sufficiently large. On the other hand if $1/2<|t|\ll 1$, then the second term is $\Oh(1)$. Hence, in both cases we have
    $$\Dist(\f,n\mapsto n^{it};N)^2 \geq \log\log{N} + \Oh(1).$$
\end{itemize}
Putting everything together, we get that \eqref{dist1.eq2} holds for $N$ sufficiently large.
\end{proof}

\begin{proof}[Proof of \cref{dist2}]
We have that
\begin{align*}
    \Dist(\f,n\mapsto n^{it};N)^2 &= \sum_{p\leq N}\frac{1-\Re\big(e(\xi/|I_N|)p^{-it}\big)}{p} \\
    &\geq \sum_{\exp((\log{N})^{\frac{2}{3}+\epsilon})\leq p\leq N} \frac{1-\Re\big(e(\xi/|I_N|))p^{-it}\big)}{p} \\
    & \geq \sum_{\exp((\log{N})^{\frac{2}{3}+\epsilon})\leq p\leq N}\frac{1}{p}-\Bigg|\sum_{\exp((\log{N})^{\frac{2}{3}+\epsilon})\leq p\leq N}\frac{e(\xi/|I_N|)p^{-it}}{p}\Bigg| \\ 
    &= \bigg(\frac{1}{3}-\epsilon\bigg)\log\log{N}-\Bigg|\sum_{\exp((\log{N})^{\frac{2}{3}+\epsilon})\leq p\leq N}\frac{1}{p^{1+it}}\Bigg|+\Oh(1),
\end{align*}
and so, it suffices to show that $$\Bigg|\sum_{\exp((\log{N})^{\frac{2}{3}+\epsilon})\leq p\leq N}\frac{1}{p^{1+it}}\Bigg|\ll 1.$$ 
The left-hand side above is equal to 
\begin{equation}\label{VK_simple}
    \bigg|\log\bigg(\zeta\bigg(1+\frac{1}{\log{N}}+it\bigg)\bigg) - \log\bigg(\zeta\bigg(1+\frac{1}{(\log{N})^{\frac{2}{3}+\epsilon}}+it\bigg)\bigg)\bigg|+\Oh(1).
\end{equation}
By the Vinogradov-Korobov zero-free region \cite{Korobov_zero-free,Vinogradov_zero-free}, there exists some absolute constant $C>0$, such that $\zeta(\sigma+it)\neq0$ in the region
$$\sigma\geq 1-\frac{C}{(\log|t|)^{\frac{2}{3}}(\log\log|t|)^{\frac{1}{3}}}, \quad |t|\geq 10,$$
and then
it follows using standard analytic methods
(as in \cite{richert}) that \eqref{VK_simple} is $\Oh(1)$.
\end{proof}

We are now ready to prove \cref{large_xi_thm1}.

\begin{proof}[Proof of \cref{large_xi_thm1}]
By \cref{Halasz_GS}, we have
$$\sum_{\substack{\xi\in I_N \\ |\xi|>A^{5/2}}}\Big|\Esub{n\in[N]} \f(n)\Big|
\ll
\sum_{\substack{\xi\in I_N \\ |\xi|>A^{5/2}}}
\exp\bigg(-\frac{M_0(\f,N)}{16}\bigg).
$$
Combining Lemmas \ref{dist1} and \ref{dist2}, we have that for $N$ sufficiently large
$$M_0(\f,N)
\geq \min\bigg\{\frac{\pi^2\xi^2}{2A^2}, \frac{\pi^2(\log\log{N})}{A}\bigg\} + \Oh(1).$$
Then we have
\begin{align*}
    \sum_{\substack{\xi\in I_N \\ |\xi|>A^{5/2}}}\Big|\Esub{n\in[N]} \f(n)\Big|
    \ll \max\Bigg\{\sum_{\xi>A^{5/2}}e^{-\frac{\pi^2\xi^2}{32A^2}}, \frac{|I_N|}{(\log{N})^\frac{\pi^2}{16A}}\Bigg\}.
\end{align*}
For $N$ sufficiently large,
\begin{multline*}
    \sum_{\xi>A^{5/2}}e^{-\frac{\pi^2\xi^2}{32A^2}} 
    \leq \int_{A^{5/2}-1}^\infty e^{-\frac{\pi^2 x^2}{32A^2}} \d x 
    = \frac{4A}{\pi} \int_{{\frac{\pi(A^{5/2}-1)}{4A}}}^\infty e^{-\frac{y^2}{2}} \d y \\
    \leq \frac{16A^2}{\pi^2(A^{5/2}-1)} \int_{{\frac{\pi(A^{5/2}-1)}{4A}}}^\infty ye^{-\frac{y^2}{2}} \d y
    \ll \exp\bigg(-\frac{\pi^2 A^3}{64}\bigg)
    \leq \exp\big(-(\log\log{N})^\frac{1}{3}\big). 
\end{multline*}
We can readily check that the second term in the maximum is smaller than the above bound, hence the result follows.
\end{proof}


\subsection{Proof of \cref{large_xi_thm2}}

We start with the following circle method estimate.

\begin{lemma}\label{circle_method}
Let $H_1\in\N$, $\Pc\subset[1,H_1]$ be a set of primes, and $(x_j)_{1\leq j\leq H_1}$, $(y_j)_{1-H_1\leq j\leq H_1-1}$ with $x_j,y_j\in\C$. 
For any $\epsilon>0$, we consider the set of major arcs 
$$\mathfrak{M}_{\epsilon,\Pc} = \Big\{\alpha\in[0,1]\colon \Big|\logE_{p\in\Pc}e(p\alpha)\Big|>\epsilon\Big\},$$
and then we have
$$\logE_{p\in\Pc}\Esub{j\in[H_1]} x_j\overline{y_{j-p}} 
\ll
\int_{\mathfrak{M}_{\epsilon,\Pc}}\Big|\Esub{j\in[H_1]} x_je(j\alpha)\Big|\d\alpha + \epsilon.$$
\end{lemma}

\begin{proof}
This is \cite[Lemma 8.1]{helf_radz} (see also \cite[Lemma 5.11]{Helf_Ubis}).
\end{proof}

We also need the following lemma concerning the Lebesgue measure of the major arcs in \cref{circle_method}.

\begin{lemma}\label{major_arcs_bound}
Let $H_1\geq 4$, $\Pc\subset[\frac{H_1}{2},H_1]$ be a set of primes. Then for any $\epsilon>0$, we have
$$|\mathfrak{M}_{\epsilon,\Pc}| \ll \frac{1}{\epsilon^4 H_1},$$
where $\mathfrak{M}_{\epsilon,\Pc}$ is as in \cref{circle_method}.
\end{lemma}

\begin{proof}
    This is \cite[Lemma 8.2]{helf_radz} (see also \cite[Lemma 5.12]{Helf_Ubis}).
\end{proof}

\begin{lemma}\label{bound_dyadic_set}
    Let $f,g\colon\Z\to\C$ be $1$-bounded functions, $H_1\geq 4$, $\Pc\subset[\frac{H_1}{2},H_1]$ be a set of primes. Then for any $\epsilon>0$, we have
    \begin{multline*}
        \logE_{n\in[N]}\logE_{p\in\Pc} f(n)\overline{g(n-p)} \\
        \ll \frac{1}{\epsilon^4}\bigg(\sup_{\alpha\in[0,1]}\logE_{n\in[N]}\Big|\Esub{j\in[H_1]}f(n+j)e(j\alpha)\Big| + \frac{\log{H_1}}{\log{N}}\bigg) + \epsilon + \frac{\log{H_1}}{\log{N}}.
    \end{multline*}
\end{lemma}

\begin{proof}
    First we have that
    $$\logE_{n\in[N]}\logE_{p\in\Pc} f(n)\overline{g(n-p)}
    = 
    \logE_{n\in[N]}\logE_{p\in\Pc}\Esub{j\in[H_1]}f(n+j)\overline{g(n+j-p)} + \Oh\bigg(\frac{\log{H_1}}{\log{N}}\bigg).$$
    We suppose that $f,g$ are both supported on $[N]$, and this will add an extra error term of order $\frac{\log{H_1}}{\log{N}}$. Then applying \cref{circle_method} for $f(n+j) = x_j$ and $g(n+j) = y_j$, $n\in[N]$, and then \cref{major_arcs_bound}, we obtain        
    \begin{multline*}
        \logE_{n\in[N]}\logE_{p\in\Pc}\Esub{j\in[H_1]}f(n+j)\overline{g(n+j-p)} \\
        \ll \frac{1}{\epsilon^4}\sup_{\alpha\in[0,1]}\logE_{n\in[N]}\Big|\Esub{j\in[H_1]}f(n+j)e(j\alpha)\Big| + \epsilon + \frac{\log{H_1}}{\log{N}}.
    \end{multline*}
    Dropping the assumption that $f$ is supported on $[N]$ adds an extra error term of order $\frac{1}{\epsilon^4}\cdot\frac{\log{H_1}}{\log{N}}$.
\end{proof}

Recall the definition of $A$, $H_0$, $H$, and $\Pc_N$ given at the beginning of \cref{reduction_section}.

\begin{proposition}\label{bound_by_local_Fourier_unif}
We have
\begin{multline*}
    \sum_{\substack{\xi\in I_N \\ |\xi|>A^{5/2}}} \logE_{n\in[N]}\Big|\logE_{p\in\Pc_N}\f(n+p)\Big|^2 \\
    \ll \frac{\exp\big(6(\log\log{N})^\frac{1}{3}\big)}{(\log\log{N})^\frac{1}{3}}\sup_{H_0\leq H'\leq H}\sup_{\alpha\in[0,1]}\logE_{n\in[N]}\Big|\Esub{j\in[H']}\f(n+j)e(j\alpha)\Big| + \frac{1}{(\log\log{N})^\frac{1}{3}}.
\end{multline*}
\end{proposition}

\begin{proof}
We start by expanding the square to get
\begin{multline}\label{expanding_square_eq1}
    \sum_{\substack{\xi\in I_N \\ |\xi|>A^{5/2}}}
    \logE_{n\in[N]}\Big|\logE_{p\in\Pc_N}\f(n+p)\Big|^2 
    = \sum_{\substack{\xi\in I_N \\ |\xi|>A^{5/2}}}
    \logE_{n\in[N]}\logE_{p\in\Pc_N}g_{\xi,N}(n)\f(n+p)  \\
    = \sum_{\substack{\xi\in I_N \\ |\xi|>A^{5/2}}}
    \logE_{n\in[N]}\logE_{p\in\Pc_N}\f(n)
    \overline{g_{\xi,N}(n-p)}
    + \Oh\bigg(\frac{A\sqrt{\log\log{N}}(\log{H})}{\log{N}}\bigg), 
\end{multline}
where we set $g_{\xi,N}(n) = \logE_{q\in\Pc_N}\f(n+q)$.
Note that using our choice of $A$ (see \eqref{I_N_def}) and $H$ (see \eqref{defn_H}), one can see that the error term in \eqref{expanding_square_eq1} is $\ll_\eta (\log{N})^{-\eta}$ for any $\eta\in(0,1)$, so it is negligible.
For the main term, we first decompose the set $\Pc_N$ into pairwise disjoint dyadic intervals $\Pc_k=[\frac{H_k}{2},H_k)\cap\P$, where $H_k = 2^{k+1}H_0$, for $k=0,\ldots,K$, where $K=\left\lceil\frac{\log(H/H_0)}{\log{2}}\right\rceil$.
For each $k$, we apply \cref{bound_dyadic_set} for $f=\f$, $g=g_{\xi,N}$, $\xi\in I_N$, $\Pc=\Pc_k$ and $\epsilon>0$ to be chosen later, to obtain
\begin{multline*}
    \logE_{n\in[N]}\logE_{p\in\Pc_k} \f(n)\overline{g_{\xi,N}(n-p)} \\
    \ll \frac{1}{\epsilon^4}\bigg(\sup_{\alpha\in[0,1]}\logE_{n\in[N]}\Big|\Esub{j\in[H_k]}\f(n+j)e(j\alpha)\Big| + \frac{\log{H}}{\log{N}}\bigg) + \epsilon.
\end{multline*}
Averaging this over all $k$ gives
\begin{align}\label{averaging_over_k}
    \Esub{k\in[0,K]} &
    \Big|\logE_{n\in[N]} \logE_{p\in\Pc_k} \f(n)\overline{g_{\xi,N}(n-p)}\Big| \notag \\
    & \ll \frac{1}{\epsilon^4}\bigg(\Esub{k\in[0,K]}\sup_{\alpha\in[0,1]}\logE_{n\in[N]}\Big|\Esub{j\in[H_k]}\f(n+j)e(j\alpha)\Big| + \frac{\log{H}}{\log{N}}\bigg) + \epsilon \notag \\
    & \leq \frac{1}{\epsilon^4}\bigg(\sup_{H_0\leq H'\leq H}\sup_{\alpha\in[0,1]}\logE_{n\in[N]}\Big|\Esub{j\in[H']}\f(n+j)e(j\alpha)\Big| + \frac{\log{H}}{\log{N}}\bigg) + \epsilon.
\end{align}

Now for any $1$-bounded function $F\colon\Z\to\C$,  
\begin{multline*}
    \Esub{k\in[0,K]}\Big|\logE_{p\in\Pc_k} F(p)\Big|
    \sim \frac{\log2}{\log(H/H_0)}\sum_{k=0}^K\frac{k\log2+\log{H_0}}{\log2}\Bigg|\sum_{p\in\Pc_k}\frac{F(p)}{p}\Bigg| \\
    \geq \frac{\log{H_0}}{\log(H/H_0)}\sum_{k=0}^K\Bigg|\sum_{p\in\Pc_k}\frac{F(p)}{p}\Bigg| 
    \geq \frac{\log{H_0}}{\log{H}}\Bigg|\sum_{p\in\Pc_N}\frac{F(p)}{p}\Bigg|,
\end{multline*}
so, using \eqref{eqn_L}, we have
\begin{align*}
    \Big|\logE_{p\in\Pc_N}F(p)\Big|
    = \frac{1}{\mathscr{L}}\Bigg|\sum_{p\in\Pc_N}\frac{F(p)}{p}\Bigg|
    & \sim \frac{1}{(\log\log{N})^\frac{1}{3}}\Bigg|\sum_{p\in\Pc_N}\frac{F(p)}{p}\Bigg| \\
    & \ll \frac{\log{H}}{(\log{H_0})(\log\log{N})^\frac{1}{3}}\,\Esub{k\in[0,K]}\Big|\logE_{p\in\Pc_k}F(p)\Big| \\
    & = \frac{\exp\big((\log\log{N})^\frac{1}{3}\big)}{(\log\log{N})^\frac{1}{3}}\,\Esub{k\in[0,K]}\Big|\logE_{p\in\Pc_k}F(p)\Big|.
\end{align*}
Applying this bound to the function $F(p) = \logE_{n\in[N]}\f(n)\overline{g_{\xi,N}(n-p)}$ and then using \eqref{averaging_over_k}, we obtain
\begin{align*}
    & \logE_{n\in[N]}\logE_{p\in\Pc_N}\f(n)\overline{g_{\xi,N}(n-p)} \\
    & \ll \frac{\exp\big((\log\log{N})^\frac{1}{3}\big)}{(\log\log{N})^\frac{1}{3}}\bigg(\frac{1}{\epsilon^4}\bigg(\sup_{H_0\leq H'\leq H}\sup_{\alpha\in[0,1]}\logE_{n\in[N]}\Big|\Esub{j\in[H']}\f(n+j)e(j\alpha)\Big| + \frac{\log{H}}{\log{N}}\bigg) + \epsilon\bigg).
\end{align*}
Summing the above over all $\xi\in I_N$ with $|\xi|>A^\frac{5}{2}$, we trivially bound the sum over $\xi$ to get an extra $|I_N|$ factor. Choosing $\epsilon = \exp(-(\log\log{N})^\frac{1}{3})|I_N|^{-1}$, we bound the left-hand side of \eqref{expanding_square_eq1} by
\begin{multline*}
    \frac{|I_N|^5 \exp\big(5(\log\log{N})^\frac{1}{3}\big)}{(\log\log{N})^\frac{1}{3}}\sup_{H_0\leq H'\leq H}\sup_{\alpha\in[0,1]}\logE_{n\in[N]}\Big|\Esub{j\in[H']}\f(n+j)e(j\alpha)\Big| \\
    + \frac{|I_N|^5 \exp\big(5(\log\log{N})^\frac{1}{3}\big)(\log{H})}{(\log\log{N})^\frac{1}{3}(\log{N})}
    + \frac{1}{(\log\log{N})^\frac{1}{3}}.
\end{multline*}
We crudely bound $|I_N|^5 \ll \exp((\log\log{N})^\frac{1}{3})$ to get the desired bound.
\end{proof}

In view of \cref{bound_by_local_Fourier_unif}, the following proposition will conclude the proof of \cref{large_xi_thm2}.

\begin{proposition}\label{F_unif_thm}
    We have
    $$\sup_{H_0\leq H'\leq H}\sup_{\alpha\in[0,1]}\sum_{\substack{\xi\in I_N \\ |\xi|>A^{5/2}}}\Esub{n\in[N]}\Big|\Esub{j\in[H']}\f(n+j)e(j\alpha)\Big|
    \ll \exp\big(-6(\log\log{N})^\frac{1}{3}\big).$$
\end{proposition}

We only need this result for logarithmic averages over $n$, but in view of \cref{Cesaro_to_log}, proving it for \Cesaro{} averages is stronger.

Given $H'\in\N$, we let 
$$Q_{H'} := \min((\log{N})^\frac{1}{125},(\log{H'})^5)$$
and given a $1$-bounded multiplicative function $f$, we define
$$M(f;N,Q_{H'}) := \inf_{\substack{\chi\:(q) \\ q\leq Q_{H'}}}\inf_{|t|\leq N}
\Dist(f, n\mapsto \chi(n)n^{it};N)^2,$$
where the first infimum is taken over all Dirichlet characters $\chi$ of modulus $q\leq Q_{H'}$.

\begin{theorem}[{\cite[Theorem 1.7]{MRT15}}]\label{MRT_thm}
Let $N\geq H'\geq 10$ and $f$ be a $1$-bounded multiplicative function. Then we have
$$
    \sup_{\alpha\in[0,1]}\frac{1}{N}\int_0^N \Big|\Esub{j\in[H']} f(x+j)e(j\alpha)\Big| \d x
    \ll \exp\bigg(-\frac{M(f;N,Q_{H'})}{20}\bigg) + \frac{\log\log{H'}}{\log{H'}} + \frac{1}{(\log{N})^\frac{1}{700}}.
$$
\end{theorem}

Let us also define
$$M^\ast(N,Q_{H'})
:= \inf_{\substack{\xi\in I_N \\ |\xi|>A^{5/2}}}\inf_{\substack{\chi\:(q); \: \chi\neq\chi_0 \\ q\leq Q_{H'}}}\inf_{|t|\leq N}
\Dist(\f, n\mapsto \chi(n)n^{it};N)^2,$$
where the second infimum is taken over all the non-principal Dirichlet characters $\chi$ of
modulus $q\leq Q_{H'}$.
The next two lemmas are the last ingredients needed to prove \cref{F_unif_thm}. As we will see, $M^\ast(N,Q_{H'})$ is large in the sense that for any $H'\geq2$, $\lim_{N\to\infty}M^\ast(N,Q_{H'}) = \infty$. In fact, we are able to prove this uniformly for all $\xi\in I_N$ and not just for $\xi$ in the range $|\xi|>A^{5/2}$.

\begin{lemma}\label{dist1 with chi}
Let $N,H'\geq 2$, $q\leq(\log{H'})^5$, $\chi$ be a non-principal Dirichlet character modulo $q$, 
$\xi\in I_N$ and $|t|\ll1$. Then we have
$$\Dist(\f,n\mapsto \chi(n)n^{it};N)^2 \gg \log\log{N}.$$
\end{lemma}

\begin{lemma}\label{dist2 with chi}
Let $N\geq 2$, $q\leq (\log{N})^\frac{1}{125}$, $\chi$ be a non-principal Dirichlet character modulo $q$, 
$\xi\in I_N$ and $10\leq |t|\leq N^B$, for some fixed $B>0$. Then for any $\epsilon>0,$ we have
$$\Dist(\f,n\mapsto \chi(n)n^{it}; N)^2 \geq \bigg(\frac{1}{3}-\epsilon\bigg)\log\log{N} + \Oh_{B,\epsilon}(1).$$
\end{lemma}

\begin{proof}[Proof of \cref{dist1 with chi}]
For $\Re(s)>1$, we define 
$$F(s,\chi)=\sum_{n\geq1}\frac{\f(n)\overline{\chi(n)}}{n^s}.$$ 
Let $s=1+\frac{1}{\log{N}}+it$, $|t|\ll1$. Then we have
\begin{equation}\label{L-f with chi}
    \log{F(s,\chi)} = \sum_{p}\frac{\f(p)\overline{\chi(p)}}{p^s} + \Oh(1) 
    = e\bigg(\frac{\xi}{|I_N|}\bigg)\log{L(s,\overline{\chi})} + \Oh(1).
\end{equation}
Using the classical asymptotic identity \cite[Eq. (2.2)]{GS_large_char_sums} as we did in the proof of \cref{dist1}, we have that
$$\exp(-\Dist(\f,n\mapsto \chi(n)n^{it};N)^2)=\frac{|F(s,\chi)|}{\log{N}} + \Oh(1),$$ 
and combing it with \eqref{L-f with chi}, we obtain that
\begin{align}\label{dist expansion with chi}
    \Dist(\f,&n\mapsto\chi(n)n^{it};N)^2
    = \log\log{N} - \log|F(s,\chi)| + \Oh(1) \notag \\
    & = \log\log{N} - \Re(\log{F(s,\chi)}) + \Oh(1) \notag \\
    & = \log\log{N} - \cos\bigg(\frac{2\pi\xi}{|I_N|}\bigg)\Re(\log{L(s,\overline{\chi})})
    -\sin\bigg(\frac{2\pi\xi}{|I_N|}\bigg)\Im(\log{L(s,\overline{\chi}})) + \Oh(1) \notag \\
    & = \log\log{N} - \cos\bigg(\frac{2\pi\xi}{|I_N|}\bigg)\log|L(s,\overline{\chi})|
    - \sin\bigg(\frac{2\pi\xi}{|I_N|}\bigg)\arg(L(s,\overline{\chi})) + \Oh(1) \notag \\
    & \geq \log\log{N} - \log|L(s,\overline{\chi})| + \Oh(1).
\end{align}
Now we use the following classical bound for Dirichlet L-functions
$$|L(s,\overline{\chi})|\ll \log(q(|t|+2)),$$
which holds uniformly for all $q>1$, all non-principal Dirichlet characters $\chi$
of modulus $q$, all $\sigma\geq1$ and all $t\in\R$ (see for example \cite[Theorem 8.18]{Tenenbaum_book}).
By the assumptions on $q$ and $t$, the above yields 
$$\log|L(s,\overline{\chi})|\ll \log\log\log{H'}.$$
Plugging this in \eqref{dist expansion with chi} yields that
$$\Dist(\f,n\mapsto \chi(n)n^{it};N)^2
\gg \log\bigg(\frac{\log{N}}{\log\log{H'}}\bigg) \gg \log\log{N}.$$
This concludes the proof.
\end{proof}

\begin{proof}[Proof of \cref{dist2 with chi}]
    We have that
\begin{align*}
    \Dist(\f,n\mapsto \chi(n)n^{it};N)^2 
    &= \sum_{p\leq N}\frac{1-\Re\big(e(\xi/|I_N|)\big)\overline{\chi(p)}p^{-it}\big)}{p} \\
    &\geq \sum_{\exp((\log{N})^{\frac{2}{3}+\epsilon})\leq p\leq N} 
    \frac{1-\Re\big(e(\xi/|I_N|)\overline{\chi(p)}p^{-it}\big)}{p} \\
    &\geq \sum_{\exp((\log{N})^{\frac{2}{3}+\epsilon})\leq p\leq N}
    \frac{1}{p}-\Bigg|\sum_{\exp((\log{N})^{\frac{2}{3}+\epsilon})\leq p\leq N}
    \frac{e(\xi/|I_N|)\overline{\chi(p)}p^{-it}}{p}\Bigg| \\
    &= \bigg(\frac{1}{3}-\epsilon\bigg)\log\log{N}- 
    \Bigg|\sum_{\exp((\log{N})^{\frac{2}{3}+\epsilon})\leq p\leq N}\frac{\overline{\chi(p)}}{p^{1+it}}\bigg|
    +\Oh(1),
\end{align*}
and so, it suffices to show that 
$$\Bigg|\sum_{\exp((\log{N})^{\frac{2}{3}+\epsilon})\leq p\leq N}\frac{\overline{\chi(p)}}{p^{1+it}}\Bigg|
\ll 1.$$ 
But this term is equal to 
\begin{equation}\label{VK}
\Big|\log\Big(L\Big(1+\frac{1}{\log{N}}+it,\chi\Big)\Big)
-\log\Big(L\Big(1+\frac{1}{(\log{N})^{\frac{2}{3}+\epsilon}}+it,\chi\Big)\Big)\Big|+\Oh(1).
\end{equation}
By the Vinogradov-Korobov zero-free region for Dirichlet L-functions there exists some constant $c>0$, such that $L(\sigma+it,\chi)\neq0$ in the region
$$\sigma\geq 1-\frac{c}{\log{q} + (\log|t|)^{\frac{2}{3}}(\log\log|t|)^{\frac{1}{3}}}, \quad |t|\geq 10,$$
(see \cite[p. 176]{Montogomery_ten}), and then
it follows using standard methods
(as in \cite{richert}) that \eqref{VK} is $\Oh(1)$.
Note that the zero-free region applies here since the the fact that $q\leq (\log{N})^\frac{1}{125}$ 
guarantees that there are no exceptional zeros (see \cite[p.6]{MRT15}).
\end{proof}

We are now ready to prove \cref{F_unif_thm}.

\begin{proof}[Proof of \cref{F_unif_thm}]
By \cref{MRT_thm}, we have that for $N$ sufficiently large,
\begin{multline*}
    \sup_{H_0\leq H'\leq H}\sup_{\alpha\in[0,1]}\sum_{\substack{\xi\in I_N \\ |\xi|>A^{5/2}}} \Esub{n\in[N]}\Big|\Esub{j\in[H']} \f(n+j)e(j\alpha)\Big| \\
    \ll \sup_{H_0\leq H'\leq H}\sum_{\substack{\xi\in I_N \\ |\xi|>A^{5/2}}} \exp\bigg(-\frac{M(\f,N,Q_{H'})}{20}\bigg) + \frac{|I_N|(\log\log{H})}{\log{H_0}} + \frac{|I_N|}{(\log{N})^\frac{1}{700}}.
\end{multline*}
We can readily check that the two last terms above are smaller than the desired bound.
Let us now bound the first term.

First we notice that
$M(\f,N,Q_{H'}) = \min\{M^\ast(\f,N,Q_{H'}),M_0(\f,N)\}$, thus, combing Lemmas
\ref{dist1}, \ref{dist2}, \ref{dist1 with chi} and \ref{dist2 with chi},
we have that for $N$ sufficiently large,
$$M(\f,N,Q) \geq \min\bigg\{\frac{\pi^2\xi^2}{2A^2}, \frac{\pi^2(\log\log{N})}{A}\bigg\},$$
hence
$$\sup_{H_0\leq H'\leq H}\sum_{\substack{\xi\in I_N \\ |\xi|>A^{5/2}}} \exp\bigg(-\frac{M(\f,N,Q_{H'})}{20}\bigg)
\ll \max\Bigg\{\sum_{\xi>A^{5/2}}e^{-\frac{\pi^2\xi^2}{40A^2}},\frac{|I_N|}{(\log{N})^\frac{\pi^2}{20A}}\Bigg\}.$$
For the first term in the maximum we have
\begin{align*}
    & \sum_{\xi>A^{5/2}}e^{-\frac{\pi^2\xi^2}{40A^2}} 
    \leq \int_{A^{5/2}-1}^\infty e^{-\frac{\pi^2 x^2}{40A^2}} \d x 
    = \frac{5A}{\pi} \int_{{\frac{\pi(A^{5/2}-1)}{5A}}}^\infty e^{-\frac{y^2}{2}} \d y \\
    & \leq \frac{25A^2}{\pi^2(A^{5/2}-1)} \int_{{\frac{\pi(A^{5/2}-1)}{5A}}}^\infty ye^{-\frac{y^2}{2}} \d y 
    \ll \exp\bigg(-\frac{\pi^2 A^3}{100}\bigg)
    \leq \exp\big(-6(\log\log{N})^\frac{1}{3}\big).
\end{align*}
Now the second term is clearly smaller, so the proof is complete.
\end{proof}




%SECTION
\section{Equivalence of Theorems~\ref{MT} and~\ref{thm_C}}
\label{Proof_Thm_C}

In this section we prove \cref{thm_C}. In fact, we show that Theorems~\ref{MT} and~\ref{thm_C} are equivalent.

For our purposes in this section, we reformulate \eqref{eqn1_thm_C} in \cref{thm_C} as follows: Given a bounded function $a\colon\N\to\C$, in light of \eqref{pk_and_logpk}, the left-hand side of \eqref{eqn1_thm_C} is 
\begin{align*}
    & \sum_{\ell=1}^\infty\overline{\pi}_\ell(N)\Big|\logE_{n\in[N]\cap\P_\ell}a(\Omega(n+1)) - \Esub{n\in[N]}a(\Omega(n))\Big| \\
    & \hspace*{0.2cm} = 
    \sum_{\ell=1}^\infty\Big|\Big(\logE_{n\in[N]\cap\P_\ell}a(\Omega(n+1))\Big)\overline{\pi}^{\log}_\ell(N) - \Big(\Esub{n\in[N]}a(\Omega(n))\Big)\overline{\pi}_\ell(N)\Big| + \Oh\Bigg(\frac{(\log\log\log{N})^2}{\sqrt{\log\log{N}}}\Bigg).
\end{align*}
Now observe that 
$$\Big(\logE_{n\in[N]\cap\P_\ell}a(\Omega(n+1))\Big)\overline{\pi}^{\log}_\ell(N)
= \logE_{n\in[N]} a(\Omega(n+1))\1_{\P_\ell}(n),$$
hence \eqref{eqn1_thm_C} is equivalent to
\begin{equation}\label{eqn_thm_C_reformulation}
    \sum_{\ell=1}^\infty \Big|\logE_{n\in[N]}a(\Omega(n+1))\1_{\P_\ell}(n) - \Big(\Esub{n\in[N]}a(\Omega(n))\Big)\overline{\pi}_\ell(N)\Big|
    \ll \frac{\sqrt{\log\log\log{N}}}{(\log\log{N})^\frac{1}{6}}.
\end{equation}

We first show the easy direction that \cref{thm_C} implies \cref{MT}.

\begin{proof}[Proof of \cref{thm_C}$\implies$\cref{MT}]
    Let $a,b\colon\N\to\C$ be bounded and suppose that \cref{thm_C} is true. Then, for $N$ sufficiently large, applying \eqref{eqn_thm_C_reformulation} to $b$, we have
    \begin{align*}
        \logE_{n\in[N]} a(\Omega(n))b(\Omega(n+1))
        & = \sum_{\ell=1}^\infty a(\ell)\Big(\logE_{n\in[N]}\1_{\P_\ell}(n)b(\Omega(n+1))\Big) \\
        & = \sum_{\ell=1}a(\ell)\Big(\logE_{n\in[N]}\1_{\P_\ell}(n)b(\Omega(n+1))\Big) + \Oh\Bigg(\frac{\sqrt{\log\log\log{N}}}{(\log\log{N})^\frac{1}{6}}\Bigg) \\
        & = \Bigg(\sum_{\ell=1}^\infty a(\ell)\overline{\pi}^{\log}_{\ell}(N)\Bigg) \Big(\logE_{n\in[N]}b(\Omega(n))\Big) + \Oh\Bigg(\frac{\sqrt{\log\log\log{N}}}{(\log\log{N})^\frac{1}{6}}\Bigg) \\
        & = \Big(\logE_{n\in[N]}a(\Omega(n))\Big) \Big(\logE_{n\in[N]}b(\Omega(n))\Big) + \Oh\Bigg(\frac{\sqrt{\log\log\log{N}}}{(\log\log{N})^\frac{1}{6}}\Bigg), 
    \end{align*}
    where the last equality follows from \eqref{eqn_expanding_a(Omega(n))_2}.
    Finally, in light of \cref{aOmega_and_logaOmega}, we can substitute both logarithmic averages in the last expression above by the corresponding \Cesaro{} averages. This concludes the proof.
\end{proof}

Finally, we show the inverse implication, thus, establishing \cref{thm_C}.

\begin{proof}[Proof of \cref{MT}$\implies$\cref{thm_C}]
Let $N\in\N$ be sufficiently large and $a\colon\N\to\C$. We can assume that $a$ is $1$-bounded.
We prove \eqref{eqn1_thm_C} and then \eqref{eqn2_thm_C} follows by Markov's inequality combined with the quantitative \Erdos-Kac theorem \cite[Theorem 2]{renyi-turan}. As we saw, it suffices to prove \eqref{eqn_thm_C_reformulation}.

By \cref{thm_A}, we have that for any $1$-bounded function $b\colon\N\to\C$,
$$\Big|\logE_{n\in[N]}a(\Omega(n+1))b(\Omega(n)) - \Big(\Esub{n\in[N]}a(\Omega(n))\Big)\Big(\Esub{n\in[N]}b(\Omega(n)\Big)\Big| 
\ll \frac{\sqrt{\log\log\log{N}}}{(\log\log{N})^\frac{1}{6}}.$$

Therefore it is enough to prove the following:
If there is $\epsilon>0$ and $N_1<N_2<\ldots\in\N$, such that for all $i\in\N$,
\begin{equation}\label{assumption_in_proof_of_thmC}
    \sum_{\ell=1}^\infty \Big|\logE_{n\in[N_i]} a(\Omega(n+1))\1_{\P_\ell}(n) - \Big(\Esub{n\in[N_i]}a(\Omega(n))\Big)\overline{\pi}_\ell(N_i)\Big|
    \geq \epsilon,
\end{equation}
then there exists some $1$-bounded function $b\colon\N\to\C$, such that for all $i\in\N$,
\begin{equation}\label{conclusion_in_proof_of_thmC}
    \Big|\logE_{n\in[N_i]}a(\Omega(n+1))b(\Omega(n)) - \Big(\Esub{n\in[N_i]}a(\Omega(n))\Big)\Big(\Esub{n\in[N_i]}b(\Omega(n)\Big)\Big| \geq \epsilon.
\end{equation}
Observe that for $\ell>N_i$, $\overline{\pi}_\ell(N_i) = 0$ and $\1_{\P_\ell}(n) = 0$ for all $n$, hence \eqref{assumption_in_proof_of_thmC} becomes
\begin{equation}
\label{assumption2_in_proof_of_thmC}
    \sum_{\ell=1}^{N_i} \Big|\logE_{n\in[N_i]} a(\Omega(n+1))\1_{\P_\ell}(n) - \Big(\Esub{n\in[N_i]}a(\Omega(n))\Big)\overline{\pi}_\ell(N_i)\Big|
    \geq \epsilon.
\end{equation}
By passing to a subsequence of $(N_i)_{i\in\N}$ if necessary, we can assume without loss of generality that
\begin{align}
\label{eqn_A_implies_first_half_of_Aprimeprimeprime_2}
\sum_{\ell=1}^{N_{i-1}}\1_{\P_\ell}(n) \leq \frac{\epsilon}{4}
\quad\text{and}\quad
\sum_{\ell=1}^{N_{i-1}} \overline{\pi}_\ell(N_i)\leq \frac{\epsilon}{4}.
\end{align}
Let $\operatorname{sgn}\colon \R\to \{-1,0,1\}$ denote the function
\[
\operatorname{sgn}(x)=
\begin{cases}
-1,&\text{if}~x<0,\\
0,&\text{if}~x=0,\\
1,&\text{if}~x>0,\\
\end{cases}
\]
and we define $b\colon\N \to \{-1,0,1\}$ as
\[
b(\ell)=\operatorname{sgn}\Big(\logE_{n\in[N_i]} a(\Omega(n+1))\1_{\P_\ell}(n) - \Big(\Esub{n\in[N_i]}a(\Omega(n))\Big)\overline{\pi}_\ell(N_i)\Big)
\]
if $\ell\in (N_{i-1},N_i]$ for some $i\geq 2$ and $b(\ell)=0$ otherwise.

First, we invoke \eqref{eqn_A_implies_first_half_of_Aprimeprimeprime_2} and the triangle inequality to deduce that
\begin{multline*}
\sum_{\ell=1}^{N_{i-1}} \Big|\logE_{n\in[N_i]}  a(\Omega(n+1))\1_{\P_\ell}(n) -\Big(\Esub{n\in[N_i]}a(\Omega(n))\Big)\overline{\pi}_\ell(N_i)\Big| \\
\leq 
\Big|\logE_{n\in[N_i]}a(\Omega(n+1))\Big|\Bigg(\sum_{\ell=1}^{N_{i-1}} \1_{\P_\ell}(n)\Bigg) + \Big|\Esub{n\in[N_i]}a(\Omega(n))\Big|\Bigg(\sum_{\ell=1}^{N_{i-1}} \overline{\pi}_\ell(N_i)\Bigg) 
\leq
\frac{\epsilon}{2}.
\end{multline*}
Combined with \eqref{assumption2_in_proof_of_thmC}, this gives
\[
\sum_{\ell=N_{i-1}+1}^{N_i} \Big|\logE_{n\in[N_i]} a(\Omega(n+1))\1_{\P_\ell}(n)
- \Big(\Esub{n\in[N_i]} a(\Omega(n))\Big)\overline{\pi}_\ell(N_i)\Big|
\geq \frac{\epsilon}{2}.
\]
Using the definition of $b(\ell)$, it now follows that 
\[
\sum_{\ell=N_{i-1}+1}^{N_i} b(\ell)\bigg(\logE_{n\in[N_i]} a(\Omega(n+1))\1_{\P_\ell}(n) 
- \Big(\Esub{n\in[N_i]} a(\Omega(n))\Big)\overline{\pi}_\ell(N_i)\bigg)
\geq \frac{\epsilon}{2}.
\]
Using \eqref{eqn_A_implies_first_half_of_Aprimeprimeprime_2} 
and the triangle inequality again we get
\begin{multline}\label{last_step_pf_thm_C}
    \Bigg|\sum_{\ell=1}^\infty b(\ell) \bigg(\logE_{n\in[N_i]} a(\Omega(n+1))\1_{\P_\ell}(n) - \Big(\Esub{n\in[N_i]} a(\Omega(n))\Big)\overline{\pi}_\ell(N_i)\bigg)\Bigg| \\
    = \Bigg|\sum_{\ell=1}^{N_i} b(\ell) \bigg(\logE_{n\in[N_i]} a(\Omega(n+1))\1_{\P_\ell}(n) - \Big(\Esub{n\in[N_i]} a(\Omega(n))\Big)\overline{\pi}_\ell(N_i)\bigg)\Bigg|
    \geq \epsilon.
\end{multline}
By rearranging the first expression above and then invoking \eqref{eqn_expanding_a(Omega(n))_1} and \eqref{eqn_expanding_a(Omega(n))_2}, 
we have 
\begin{align*}
    \sum_{\ell=1}^\infty b(\ell) & \bigg(\logE_{n\in[N_i]} a(\Omega(n+1))\1_{\P_\ell}(n) - \Big(\Esub{n\in[N_i]} a(\Omega(n))\Big)\overline{\pi}_\ell(N_i)\bigg) \\
    & = \logE_{n\in[N_i]}a(\Omega(n+1))\Bigg(\sum_{\ell=1}^\infty b(\ell)\1_{\P_\ell}(n)\Bigg) - \Big(\Esub{n\in[N_i]}a(\Omega(n))\Big)\Bigg(\sum_{\ell=1}^\infty b(\ell)\overline{\pi}_\ell(N_i)\Bigg) \\
    & = \logE_{n\in[N_i]}a(\Omega(n+1))b(\Omega(n)) - \Big(\Esub{n\in[N_i]}a(\Omega(n))\Big)\Big(\Esub{n\in[N_i]}b(\Omega(n))\Big).
\end{align*}
Combining this with \eqref{last_step_pf_thm_C}, \eqref{conclusion_in_proof_of_thmC} follows, thus, concluding the proof.
\end{proof}

%APPENDIX
%\appendix
%\section{Appendix}

%BIBLIOGRAPHY
%\bibliographystyle{plain}
%\bibliographystyle{aomplain}
\bibliographystyle{aomalpha}
\bibliography{Refs_Omega.bib}

%\printbibliography
%\begin{thebibliography}{9}
%\bibitem{}
%\end{thebibliography} 


%==========================================================
%==========================================================

%AFFILIATION
\bigskip
\bigskip
\footnotesize
\noindent
Dimitrios Charamaras\\
\textsc{{\'E}cole Polytechnique F{\'e}d{\'e}rale de Lausanne (EPFL)}\\
\href{mailto:dimitrios.charamaras@epfl.ch}
{\texttt{dimitrios.charamaras@epfl.ch}}

\bigskip
\noindent
Florian K. Richter\\
\textsc{{\'E}cole Polytechnique F{\'e}d{\'e}rale de Lausanne (EPFL)}\\
\href{mailto:f.richter@epfl.ch}
{\texttt{f.richter@epfl.ch}}

%END DOCUMENT
\end{document}
\providecommand{\etalchar}[1]{$^{#1}$}
\providecommand{\bysame}{\leavevmode\hbox to3em{\hrulefill}\thinspace}
\providecommand{\noopsort}[1]{}
\providecommand{\mr}[1]{\href{http://www.ams.org/mathscinet-getitem?mr=#1}{MR~#1}}
\providecommand{\zbl}[1]{\href{http://www.zentralblatt-math.org/zmath/en/search/?q=an:#1}{Zbl~#1}}
\providecommand{\jfm}[1]{\href{http://www.emis.de/cgi-bin/JFM-item?#1}{JFM~#1}}
\providecommand{\arxiv}[1]{\href{http://www.arxiv.org/abs/#1}{arXiv~#1}}
\providecommand{\doi}[1]{\url{https://doi.org/#1}}
\providecommand{\MR}{\relax\ifhmode\unskip\space\fi MR }
% \MRhref is called by the amsart/book/proc definition of \MR.
\providecommand{\MRhref}[2]{%
  \href{http://www.ams.org/mathscinet-getitem?mr=#1}{#2}
}
\providecommand{\href}[2]{#2}
\begin{thebibliography}{EAKPLdlR17}

\bibitem[BR22]{BR22}
\bgroup\scshape{}V.~Bergelson\egroup{} and \bgroup\scshape{}F.~K. Richter\egroup{}, {D}ynamical generalizations of the prime number theorem and disjointness of additive and multiplicative semigroup actions,  \emph{Duke Math. J.} \textbf{171} no.~15 (2022), 3133--3200.

\bibitem[Cha25]{Charamaras_MVT_24}
\bgroup\scshape{}D.~Charamaras\egroup{}, Mean value theorems in multiplicative systems and joint ergodicity of additive and multiplicative actions,  \emph{Trans. Amer. Math. Soc.} \textbf{378} no.~3 (2025), 1883–1937.

\bibitem[Cho87]{Chowla87}
\bgroup\scshape{}S.~Chowla\egroup{}, \emph{The {R}iemann {H}ypothesis and {H}ilbert's {T}enth {P}roblem}, CRC Press, 1987.

\bibitem[EAKPLdlR17]{EAKLdlR17}
\bgroup\scshape{}E.~H. El~Abdalaoui\egroup{}, \bgroup\scshape{}J.~Ku{\l}aga-Przymus\egroup{}, \bgroup\scshape{}M.~Lema{\'n}czyk\egroup{}, and \bgroup\scshape{}T.~de~la Rue\egroup{}, The {C}howla and the {S}arnak conjectures from ergodic theory point of view,  \emph{Discrete Contin. Dyn. Syst.} \textbf{37} no.~6 (2017), 2899--2944.

\bibitem[Ell77]{elliott_TK}
\bgroup\scshape{}P.~D. T.~A. Elliott\egroup{}, The {T}ur\'an-{K}ubilius inequality,  \emph{Proc. Amer. Math. Soc.} \textbf{65} no.~1 (1977), 8--10.

\bibitem[Ell92]{Elliott92}
\bgroup\scshape{}P.~D. T.~A. Elliott\egroup{}, On the correlation of multiplicative functions,  \emph{Notas Soc. Mat. Chile} \textbf{11} no.~1 (1992), 1--11.

\bibitem[Erd48]{Erdos48a}
\bgroup\scshape{}P.~Erd\H{o}s\egroup{}, On the integers having exactly {$K$} prime factors,  \emph{Ann. of Math. (2)} \textbf{49} (1948), 53--66.

\bibitem[FKPL18]{FKL18}
\bgroup\scshape{}S.~Ferenczi\egroup{}, \bgroup\scshape{}J.~Ku{\l}aga-Przymus\egroup{}, and \bgroup\scshape{}M.~Lema{\'n}czyk\egroup{}, Sarnak's conjecture: what's new,  in \emph{Ergodic theory and dynamical systems in their interactions with arithmetics and combinatorics}, \emph{Lecture Notes in Math.} \textbf{2213}, Springer, Cham, 2018, pp.~163--235.

\bibitem[Fra17]{Frantzikinakis17}
\bgroup\scshape{}N.~Frantzikinakis\egroup{}, Ergodicity of the {L}iouville system implies the {C}howla conjecture,  \emph{Discrete Anal.} (2017), Paper No. 19, 41.

\bibitem[FH18]{FH18a}
\bgroup\scshape{}N.~Frantzikinakis\egroup{} and \bgroup\scshape{}B.~Host\egroup{}, The logarithmic {S}arnak conjecture for ergodic weights,  \emph{Ann. of Math. (2)} \textbf{187} no.~3 (2018), 869--931.

\bibitem[GLdlR21]{GLdlR21}
\bgroup\scshape{}A.~Gomilko\egroup{}, \bgroup\scshape{}M.~Lema{\'n}czyk\egroup{}, and \bgroup\scshape{}T.~de~la Rue\egroup{}, On {F}urstenberg systems of aperiodic multiplicative functions of {M}atom\"{a}ki, {R}adziwi{\l}{\l} and {T}ao,  \emph{J. Mod. Dyn.} \textbf{17} (2021), 529--555.

\bibitem[GS03]{granville2003decay}
\bgroup\scshape{}A.~Granville\egroup{} and \bgroup\scshape{}K.~Soundararajan\egroup{}, Decay of mean values of multiplicative functions,  \emph{Canad. J. Math.} \textbf{55} no.~6 (2003), 1191--1230.

\bibitem[GS18]{GS_large_char_sums}
\bgroup\scshape{}A.~Granville\egroup{} and \bgroup\scshape{}K.~Soundararajan\egroup{}, Large character sums: {B}urgess's theorem and zeros of {$L$}-functions,  \emph{J. Eur. Math. Soc. (JEMS)} \textbf{20} no.~1 (2018), 1--14.

\bibitem[HU23]{Helf_Ubis}
\bgroup\scshape{}H.~Helfgott\egroup{} and \bgroup\scshape{}A.~Ubis\egroup{}, Primes, parity and analysis,  \emph{Publ. Mat. Urug.} \textbf{18} (2023), 205--283.

\bibitem[HR21]{helf_radz}
\bgroup\scshape{}H.~A. Helfgott\egroup{} and \bgroup\scshape{}M.~Radziwi{\l}{\l}\egroup{}, Expansion, divisibility and parity,  \emph{arXiv preprint arXiv:2103.06853} (2021).

\bibitem[KLRT24]{KLRT2024}
\bgroup\scshape{}A.~Kanigowski\egroup{}, \bgroup\scshape{}M.~Lema\'nczyk\egroup{}, \bgroup\scshape{}F.~K. Richter\egroup{}, and \bgroup\scshape{}J.~Ter\"av\"ainen\egroup{}, On the local {F}ourier uniformity problem for small sets,  \emph{Int. Math. Res. Not. IMRN} no.~15 (2024), 11488--11512.

\bibitem[KMT23]{KMT23_almost_all_scales}
\bgroup\scshape{}O.~Klurman\egroup{}, \bgroup\scshape{}A.~P. Mangerel\egroup{}, and \bgroup\scshape{}J.~Ter{\"a}v{\"a}inen\egroup{}, On {E}lliott's conjecture and applications,  \emph{arXiv preprint arXiv:2304.05344} (2023).

\bibitem[Kor58]{Korobov_zero-free}
\bgroup\scshape{}N.~M. Korobov\egroup{}, Estimates of trigonometric sums and their applications,  \emph{Uspehi Mat. Nauk} \textbf{13} no.~4(82) (1958), 185--192.

\bibitem[Loy23]{Loyd23}
\bgroup\scshape{}K.~Loyd\egroup{}, {A} dynamical approach to the asymptotic behavior of the sequence {$\Omega(n)$},  \emph{Ergodic Theory Dynam. Systems} \textbf{43} no.~11 (2023), 3685--3706.

\bibitem[MRT15]{MRT15}
\bgroup\scshape{}K.~Matom\"{a}ki\egroup{}, \bgroup\scshape{}M.~Radziwi{\l}{\l}\egroup{}, and \bgroup\scshape{}T.~Tao\egroup{}, An averaged form of {C}howla's conjecture,  \emph{Algebra Number Theory} \textbf{9} no.~9 (2015), 2167--2196.

\bibitem[MRS{\etalchar{+}}24]{MRSTT2024}
\bgroup\scshape{}K.~Matom{\"a}ki\egroup{}, \bgroup\scshape{}M.~Radziwi{\l}{\l}\egroup{}, \bgroup\scshape{}X.~Shao\egroup{}, \bgroup\scshape{}T.~Tao\egroup{}, and \bgroup\scshape{}J.~Ter{\"a}v{\"a}inen\egroup{}, Higher uniformity of arithmetic functions in short intervals {II}. {A}lmost all intervals,  \emph{arXiv preprint arXiv:2411.05770} (2024).

\bibitem[MRT{\etalchar{+}}23]{MRTT2023_higher_uniformity}
\bgroup\scshape{}K.~Matom\"aki\egroup{}, \bgroup\scshape{}M.~Radziwi\l\l\egroup{}, \bgroup\scshape{}T.~Tao\egroup{}, \bgroup\scshape{}J.~Ter\"av\"ainen\egroup{}, and \bgroup\scshape{}T.~Ziegler\egroup{}, Higher uniformity of bounded multiplicative functions in short intervals on average,  \emph{Ann. of Math. (2)} \textbf{197} no.~2 (2023), 739--857.

\bibitem[MSTT23]{MSTT2023}
\bgroup\scshape{}K.~Matom\"aki\egroup{}, \bgroup\scshape{}X.~Shao\egroup{}, \bgroup\scshape{}T.~Tao\egroup{}, and \bgroup\scshape{}J.~Ter\"av\"ainen\egroup{}, Higher uniformity of arithmetic functions in short intervals {I}. {A}ll intervals,  \emph{Forum Math. Pi} \textbf{11} (2023), Paper No. e29, 97.

\bibitem[Mon94]{Montogomery_ten}
\bgroup\scshape{}H.~L. Montgomery\egroup{}, \emph{Ten lectures on the interface between analytic number theory and harmonic analysis}, \emph{CBMS Regional Conference Series in Mathematics} \textbf{84}, Conference Board of the Mathematical Sciences, Washington, DC; by the American Mathematical Society, Providence, RI, 1994.

\bibitem[MV07]{MV_book}
\bgroup\scshape{}H.~L. Montgomery\egroup{} and \bgroup\scshape{}R.~C. Vaughan\egroup{}, \emph{Multiplicative number theory. {I}. {C}lassical theory}, \emph{Cambridge Studies in Advanced Mathematics} \textbf{97}, Cambridge University Press, Cambridge, 2007.

\bibitem[Pil23]{pilatte2023improved}
\bgroup\scshape{}C.~Pilatte\egroup{}, Improved bounds for the two-point logarithmic {C}howla conjecture,  \emph{arXiv preprint arXiv:2310.19357} (2023).

\bibitem[Ram18]{Ramare18}
\bgroup\scshape{}O.~Ramar\'{e}\egroup{}, Chowla's conjecture: from the {L}iouville function to the {M}oebius function,  in \emph{Ergodic theory and dynamical systems in their interactions with arithmetics and combinatorics}, \emph{Lecture Notes in Math.} \textbf{2213}, Springer, Cham, 2018, pp.~317--323.

\bibitem[RT58]{renyi-turan}
\bgroup\scshape{}A.~R\'{e}nyi\egroup{} and \bgroup\scshape{}P.~Tur\'{a}n\egroup{}, On a theorem of {E}rd{\H{o}}s-{K}ac,  \emph{Acta Arith.} \textbf{4} (1958), 71--84.

\bibitem[Ric67]{richert}
\bgroup\scshape{}H.-E. Richert\egroup{}, Zur {A}bsch\"atzung der {R}iemannschen {Z}etafunktion in der {N}\"ahe der {V}ertikalen {$\sigma =1$},  \emph{Math. Ann.} \textbf{169} (1967), 97--101.

\bibitem[Ric21]{Richter_PNT_21}
\bgroup\scshape{}F.~K. Richter\egroup{}, A new elementary proof of the {P}rime {N}umber {T}heorem,  \emph{Bull. Lond. Math. Soc.} \textbf{53} no.~5 (2021), 1365--1375.

\bibitem[Sar11]{Sarnak11}
\bgroup\scshape{}P.~Sarnak\egroup{}, \emph{Three {L}ectures on the {M}obius function randomness and dynamics}, 2011. Available at \url{publications.ias.edu/sarnak}.

\bibitem[Sel54]{Selberge_Sathe}
\bgroup\scshape{}A.~Selberg\egroup{}, Note on a paper by {L}. {G}. {S}athe,  \emph{J. Indian Math. Soc. (N.S.)} \textbf{18} (1954), 83--87.

\bibitem[Tao16]{Tao16}
\bgroup\scshape{}T.~Tao\egroup{}, The logarithmically averaged {C}howla and {E}lliott conjectures for two-point correlations,  \emph{Forum Math. Pi} \textbf{4} (2016), e8, 36.

\bibitem[TT18]{TT18}
\bgroup\scshape{}T.~Tao\egroup{} and \bgroup\scshape{}J.~Ter\"{a}v\"{a}inen\egroup{}, Odd order cases of the logarithmically averaged {C}howla conjecture,  \emph{J. Th\'{e}or. Nombres Bordeaux} \textbf{30} no.~3 (2018), 997--1015.

\bibitem[Tao17]{tao2017equivalence}
\bgroup\scshape{}T.~Tao\egroup{}, Equivalence of the logarithmically averaged {C}howla and {S}arnak conjectures,  in \emph{Number theory---{D}iophantine problems, uniform distribution and applications}, Springer, Cham, 2017, pp.~391--421.

\bibitem[Ten15]{Tenenbaum_book}
\bgroup\scshape{}G.~Tenenbaum\egroup{}, \emph{Introduction to analytic and probabilistic number theory}, third ed., \emph{Graduate Studies in Mathematics} \textbf{\textbf{163}}, American Mathematical Society, Providence, RI, 2015.

\bibitem[Vin58]{Vinogradov_zero-free}
\bgroup\scshape{}I.~M. Vinogradov\egroup{}, A new estimate of the function {$\zeta (1+it)$},  \emph{Izv. Akad. Nauk SSSR Ser. Mat.} \textbf{22} (1958), 161--164.

\end{thebibliography}
